\documentclass[12pt,a4paper]{report}

% =========================================================
% CORE PACKAGES & GEOMETRY
% =========================================================
\usepackage[utf8]{inputenc}
\usepackage[T1]{fontenc}
\usepackage{newtxtext,newtxmath}
\usepackage[english]{babel}
\usepackage{microtype}
\usepackage{graphicx}
\usepackage{geometry}
\geometry{a4paper, top=1.0in, bottom=1.0in, left=1.25in, right=1.0in}
\usepackage{amsmath}
% Preemptively undefine \Bbbk to avoid conflict between newtxmath and any AMS font load
\let\Bbbk\relax
\usepackage{cite}
\usepackage{setspace}
\usepackage{booktabs}
\usepackage{tabularx}
\usepackage{caption}
\usepackage{float}
\usepackage{fancyhdr}
\usepackage{adjustbox}
\usepackage{enumitem}
\usepackage{titlesec}
\usepackage{tocloft}
\usepackage{needspace}
\usepackage{tikz}
\usetikzlibrary{shapes.geometric, shapes.multipart, arrows.meta, calc, positioning, backgrounds}
\usepackage{hyperref}

% ---------------------------------------------------------
% TYPOGRAPHY & OVERFLOW ELIMINATION
% ---------------------------------------------------------
\setstretch{1.45}
\setlength{\parindent}{1.5em}
\setlength{\parskip}{0pt}  % Prevents artificial elastic whitespace
\raggedbottom              % Enforces clean bottom page alignments

\clubpenalty=10000
\widowpenalty=10000
\displaywidowpenalty=10000

\setlist{itemsep=2pt, topsep=3pt, parsep=0pt, partopsep=0pt, leftmargin=2em}
\renewcommand{\arraystretch}{1.20}
\setlength{\textfloatsep}{14pt plus 2pt minus 2pt}
\setlength{\intextsep}{12pt plus 2pt minus 2pt}

\captionsetup[figure]{justification=centering, position=below, skip=8pt, font=small, labelfont=bf}
\captionsetup[table]{justification=centering, position=above, skip=8pt, font=small, labelfont=bf}

\hypersetup{colorlinks=true, linkcolor=black, citecolor=black, urlcolor=black}
\makeatletter
\providecommand*{\toclevel@front}{0}
\makeatother

% ---------------------------------------------------------
% ORPHAN-PROOF SECTION HEADINGS
% ---------------------------------------------------------
\titleformat{\chapter}[display]{\normalfont\huge\bfseries}{\chaptertitlename\ \thechapter}{12pt}{\Huge}
\titleformat{name=\chapter,numberless}[block]{\normalfont\huge\bfseries}{}{0pt}{\Huge}
\titlespacing*{\chapter}{0pt}{-10pt}{20pt}

\titleformat{\section}{\Needspace{7\baselineskip}\normalfont\Large\bfseries}{\thesection}{0.8em}{}
\titleformat{\subsection}{\Needspace{5\baselineskip}\normalfont\large\bfseries}{\thesubsection}{0.8em}{}
\titleformat{\subsubsection}{\Needspace{4\baselineskip}\normalfont\normalsize\bfseries}{\thesubsubsection}{0.8em}{}

\titlespacing*{\section}{0pt}{14pt plus 2pt minus 2pt}{6pt plus 1pt}
\titlespacing*{\subsection}{0pt}{12pt plus 2pt minus 2pt}{4pt plus 1pt}
\titlespacing*{\subsubsection}{0pt}{10pt plus 2pt minus 2pt}{3pt plus 1pt}

% ---------------------------------------------------------
% TOC, LOF, LOT STYLING
% ---------------------------------------------------------
\renewcommand{\contentsname}{INDEX}
\renewcommand{\listfigurename}{List of Figures}
\renewcommand{\listtablename}{List of Tables}
\setcounter{tocdepth}{3}
\setcounter{secnumdepth}{3}

\renewcommand{\cfttoctitlefont}{\hfill\Huge\bfseries}
\renewcommand{\cftaftertoctitle}{\hfill}
\renewcommand{\cftloftitlefont}{\Huge\bfseries}
\renewcommand{\cftlottitlefont}{\Huge\bfseries}
\renewcommand{\cftdotsep}{2}
\renewcommand{\cftchapfont}{\bfseries}
\renewcommand{\cftchappagefont}{\bfseries}
\renewcommand{\cftchapleader}{\dotfill}
\setlength{\cftbeforechapskip}{8pt}
\setlength{\cftbeforesecskip}{2pt}
\setlength{\cftbeforesubsecskip}{2pt}
\setlength{\cftchapindent}{0em}   \setlength{\cftchapnumwidth}{2.0em}
\setlength{\cftsecindent}{2.0em}  \setlength{\cftsecnumwidth}{2.6em}
\setlength{\cftsubsecindent}{4.6em} \setlength{\cftsubsecnumwidth}{3.4em}
\setlength{\cftfigindent}{1.5em}  \setlength{\cftfignumwidth}{2.8em}
\setlength{\cfttabindent}{1.5em}  \setlength{\cfttabnumwidth}{2.8em}

\makeatletter
\newcommand*{\l@front}[2]{\@dottedtocline{0}{0em}{0em}{#1}{#2}}
\makeatother

% ---------------------------------------------------------
% RUNNING HEADERS & FOOTERS
% ---------------------------------------------------------
\setlength{\headheight}{15pt}
\pagestyle{fancy}
\fancyhf{}
\renewcommand{\headrulewidth}{0.4pt}
\renewcommand{\footrulewidth}{0.4pt}
\fancyhead[L]{\small Resource-Aware Collaborative Cloud IDE}
\fancyhead[R]{\small Stage-I Project Report}
\fancyfoot[C]{\small Department of Computer Engineering, AIT Pune}
\fancyfoot[R]{\small\thepage}

\fancypagestyle{plain}{%
  \fancyhf{}
  \renewcommand{\headrulewidth}{0pt}
  \renewcommand{\footrulewidth}{0.4pt}
  \fancyfoot[C]{\small Department of Computer Engineering, AIT Pune}
  \fancyfoot[R]{\small\thepage}
}

\newcommand{\pageframe}{%
  \begin{tikzpicture}[remember picture,overlay]
    \draw[line width=0.8pt] ($(current page.north west)+(1.2cm,-1.2cm)$) rectangle ($(current page.south east)+(-1.2cm,1.2cm)$);
  \end{tikzpicture}}

% ---------------------------------------------------------
% GRAPHICS & FALLBACK HELPERS
% ---------------------------------------------------------
\newcommand{\safeincludegraphics}[3][width=0.85\linewidth]{%
  \IfFileExists{#2}{%
    \includegraphics[#1]{#2}%
  }{%
    \fbox{\parbox{0.8\linewidth}{\centering\textbf{#3}\\[0.2cm]{\small [Asset Placeholder: \texttt{\detokenize{#2}}]}}}%
  }%
}

\newcommand{\pngonly}[1]{%
  \IfFileExists{#1}{%
    \adjustbox{max width=\linewidth, max height=0.70\textheight, keepaspectratio, center}{\includegraphics{#1}}%
  }{\fbox{\parbox{0.8\linewidth}{\centering Missing Image Asset: \texttt{\detokenize{#1}}}}}}

\newcommand{\pngdiag}[2]{%
  \IfFileExists{#1}{%
    \adjustbox{max width=\linewidth, max height=0.70\textheight, keepaspectratio, center}{\includegraphics{#1}}%
  }{\adjustbox{max width=\linewidth, max height=0.70\textheight, keepaspectratio, center}{#2}}}

\newcommand{\pngdiagrot}[2]{%
  \IfFileExists{#1}{%
    \adjustbox{angle=90, max width=\linewidth, max height=0.70\textheight, keepaspectratio, center}{\includegraphics{#1}}%
  }{\adjustbox{angle=90, max width=\linewidth, max height=0.70\textheight, keepaspectratio, center}{#2}}}

\newcommand{\gl}[1]{\guillemotleft\,#1\,\guillemotright}
\newcommand{\LT}{\textless}
\newcommand{\GT}{\textgreater}

\tikzset{
  >={Stealth[length=2.4mm,width=2mm]},
  lbl/.style={font=\footnotesize, inner sep=1.2pt, fill=white},
  tl/.style={font=\footnotesize, inner sep=1.2pt, fill=white, align=center},
  sobj/.style={rectangle, draw=black, thick, fill=white, text=black, minimum width=3.4cm, minimum height=1.2cm,
               align=center, font=\small\bfseries, inner sep=3pt},
  sm/.style={-{Straight Barb[length=2.6mm,width=2.8mm]}, thick, black},
  sr/.style={-{Straight Barb[length=2.6mm,width=2.8mm]}, thick, dashed, black}
}

\newcommand{\actor}[4]{%
  \draw[thick] (#2,#3+0.55) circle (0.17);
  \draw[thick] (#2,#3+0.38)--(#2,#3-0.2) (#2-0.35,#3+0.15)--(#2+0.35,#3+0.15)
        (#2,#3-0.2)--(#2-0.28,#3-0.8) (#2,#3-0.2)--(#2+0.28,#3-0.8);
  \node[font=\footnotesize\bfseries,align=center,text width=2.8cm] at (#2,#3-1.5) {#4};
  \coordinate (#1e) at (#2+0.38,#3+0.15);
  \coordinate (#1w) at (#2-0.38,#3+0.15);}

\newcommand{\sq}[5]{%
  \pgfmathsetmacro{\sqW}{max(abs(#2-#3)-0.5,1.5)}
  \draw[#1] (#2,#4) -- (#3,#4);
  \node[font=\footnotesize, text width=\sqW cm, align=center, anchor=south, inner sep=1.5pt]
     at ({(#2+#3)/2},#4) {#5};}
\newcommand{\sql}[6]{%
  \draw[#1] (#2,#4) -- (#3,#4);
  \node[font=\footnotesize, fill=white, anchor=south, inner sep=1.5pt, yshift=1pt]
     at ({#2+(#3-#2)*#6},#4) {#5};}
\newcommand{\selfm}[4]{%
  \draw[thick,black,->] (#1+0.1,#2) -- (#1+0.75,#2) -- (#1+0.75,#2-0.5) -- (#1+0.1,#2-0.5);
  \node[anchor=west,font=\footnotesize,text width=#4 cm,align=left,inner sep=1.5pt]
     at (#1+0.8,#2-0.25) {#3};}
\newcommand{\act}[3]{\draw[fill=white,draw=black,thick] (#1-0.1,#2) rectangle (#1+0.1,#3);}
\newcommand{\frag}[7]{%
  \draw[thick,black] (#1,#2) rectangle (#3,#4);
  \node[anchor=north west, draw=black, fill=white, font=\footnotesize\bfseries, inner sep=2pt,
        minimum width=1.1cm] at (#1,#2) {#5};
  \node[anchor=west, font=\footnotesize, inner sep=1pt] at (#7,{#2-0.28}) {#6};}

\newcommand{\state}[5][white]{\node[st,fill=white,draw=black] (#2) at (#3) {\textbf{#4}\nodepart{two}\footnotesize #5};}
\newcommand{\inst}[4]{\node[obj] (#1) at (#2) {{\centering\underline{\textbf{\large #3}}\par}\nodepart{two}#4};}

\newcommand{\archlayer}[4]{%
  \draw[rounded corners=4pt,thick,dashed,draw=black,fill=white] (-6.6,#2) rectangle (6.6,#3);
  \node[font=\small\bfseries,anchor=north,text=black] at (0,#2-0.08) {#4};}
\newcommand{\archbox}[6]{%
  \node[rounded corners=2pt,draw=black,fill=white,thick,font=\scriptsize,align=center,
        inner sep=2pt,minimum width=#4 cm,minimum height=#5 cm,text width=#4 cm] at (#2,#3) {#6};}

\begin{document}

% =========================================================
% TITLE PAGE
% =========================================================
\begin{titlepage}
    \thispagestyle{empty}
    \pageframe
    \begin{center}
        \vspace*{0.2cm}
        {\large A} \\[0.2cm]
        {\large Preliminary Project Report On} \\[0.4cm]
        {\LARGE \textbf{Resource-Aware Collaborative Cloud IDE}} \\[0.9cm]

        {\normalsize SUBMITTED TOWARDS THE} \\
        {\normalsize PARTIAL FULFILLMENT OF THE REQUIREMENTS OF} \\[0.5cm]
        {\large \textbf{Bachelor Of Engineering (Computer Engineering)}} \\[0.7cm]

        {\large \textbf{By}} \\[0.3cm]
        \begin{tabular}{p{5.8cm}l}
            Abhay Joshi & Exam No: 7402 \\
            Anish Kumar Kaler & Exam No: 7411 \\
            Anish Kumar Pandey & Exam No: 7412 \\
            Bikash Sharma & Exam No: 7418 \\
        \end{tabular} \\[0.7cm]

        {\large \textbf{Under The Guidance of}} \\[0.2cm]
        Prof. Yogita Hambir \\[0.6cm]

        \vfill
        \safeincludegraphics[width=0.32\textwidth]{logo.png}{College Logo (logo.png)} \\[0.4cm]

        Department of Computer Engineering \\
        Army Institute of Technology, Pune-411015 \\[0.2cm]
        SAVITRIBAI PHULE PUNE UNIVERSITY \\[0.1cm]
        2026-27 \\[0.4cm]
    \end{center}
\end{titlepage}

% =========================================================
% CERTIFICATE
% =========================================================
\clearpage
\thispagestyle{empty}
\pageframe
\begin{center}
    \vspace*{0.2cm}
    \safeincludegraphics[width=0.24\textwidth]{logo.png}{College Logo} \\[0.3cm]
    {\large \textbf{ARMY INSTITUTE OF TECHNOLOGY, PUNE,}} \\
    {\large \textbf{DEPARTMENT OF COMPUTER ENGINEERING}} \\[0.35cm]
    {\Large \textbf{CERTIFICATE}} \\[0.4cm]
    This is to certify that the Project Entitled \\[0.25cm]
    {\large \textbf{Resource-Aware Collaborative Cloud IDE}} \\[0.25cm]
    Submitted by \\[0.2cm]
    \begin{tabular}{p{5.8cm}l}
        Abhay Joshi & Exam No: 7402 \\
        Anish Kumar Kaler & Exam No: 7411 \\
        Anish Kumar Pandey & Exam No: 7412 \\
        Bikash Sharma & Exam No: 7418 \\
    \end{tabular}
\end{center}
\vspace{0.25cm}
is a bonafide work carried out by students under the supervision of \mbox{Prof. Yogita Hambir} and is submitted towards the partial fulfillment of the requirements for the Bachelor of Engineering (Computer Engineering) Project.

\vfill
\noindent
\begin{minipage}[t]{0.48\textwidth}
    Prof. Yogita Hambir \\
    Internal Guide \\
    Department of Computer Engg.
\end{minipage}
\hfill
\begin{minipage}[t]{0.48\textwidth}
    \raggedleft
    Prof. (Dr.) Sunil Dhore \\
    H.O.D. \\
    Department of Computer Engg.
\end{minipage}

\vspace{1.3cm}
\noindent
\begin{minipage}[t]{0.48\textwidth}
    \rule{4.5cm}{0.4pt} \\
    External Examiner
\end{minipage}
\hfill
\begin{minipage}[t]{0.48\textwidth}
    \raggedleft
    Dr. B. P. Patil \\
    Principal
\end{minipage}

\vspace{0.5cm}
\noindent Place : AIT, Pune \hfill Date : \rule{3.5cm}{0.4pt}

% =========================================================
% PROJECT APPROVAL SHEET
% =========================================================
\clearpage
\thispagestyle{empty}
\pageframe
\begin{center}
    \vspace*{0.2cm}
    {\large \textbf{PROJECT APPROVAL SHEET}} \\[0.4cm]
    A \\[0.2cm]
    Project Stage-I Report on \\[0.3cm]
    {\large \textbf{Resource-Aware Collaborative Cloud IDE}} \\[0.35cm]
    is successfully completed by \\[0.3cm]
    \begin{tabular}{p{5.8cm}l}
        Abhay Joshi & Exam No: 7402 \\
        Anish Kumar Kaler & Exam No: 7411 \\
        Anish Kumar Pandey & Exam No: 7412 \\
        Bikash Sharma & Exam No: 7418 \\
    \end{tabular} \\[0.4cm]
    at \\[0.3cm]
    \safeincludegraphics[width=0.24\textwidth]{logo.png}{College Logo} \\[0.35cm]
    Department of Computer Engineering \\
    Army Institute of Technology, Pune-411015 \\[0.2cm]
    SAVITRIBAI PHULE PUNE UNIVERSITY \\[0.1cm]
    2026-27
\end{center}

\vfill
\noindent
\begin{minipage}[t]{0.48\textwidth}
    Prof. Yogita Hambir \\
    Internal Guide \\
    Department of Computer Engg.
\end{minipage}
\hfill
\begin{minipage}[t]{0.48\textwidth}
    \raggedleft
    Prof. (Dr.) Sunil Dhore \\
    H.O.D. \\
    Department of Computer Engg.
\end{minipage}

% =========================================================
% PRELIMINARY PAGES
% =========================================================
\clearpage
\pagenumbering{roman}
\setcounter{page}{1}

\begin{center}
    \vspace*{0.2cm}
    {\Huge \textbf{Abstract}}
\end{center}
\vspace{0.3cm}
\thispagestyle{plain}
\addcontentsline{toc}{front}{Abstract}
Cloud-based Integrated Development Environments (Cloud IDEs) have become fundamental to modern software engineering workflows. By providing pre-configured runtime environments, centralized access governance, and browser-accessible development sandboxes, Cloud IDEs eliminate local environment drift and reduce developer onboarding time. However, existing commercial and open-source implementations exhibit structural resource inefficiencies. Current architectures statically assign dedicated CPU cores and RAM limits to containerized sessions. These allocations are sized to withstand peak compilation or test-suite bursts, which accounts for less than 10\% of total session duration. During typing, reading, and conceptualization, actual container utilization plummets below 5\%, resulting in wasted cloud capacity and high server energy draw. Simultaneously, centralized collaborative operational transformation models create synchronization bottlenecks and lock contention when scaling to higher user concurrencies.

This project presents the design, mathematical modeling, and empirical realization of a \textbf{Resource-Aware Collaborative Cloud IDE}. The platform establishes an adaptive, closed-loop control system that couples developer interaction states with infrastructure-level resource management. The client tier utilizes decentralized Conflict-free Replicated Data Types (CRDTs) built on Yjs, ensuring mathematically verified eventual consistency, sub-100ms peer synchronization latency, and zero central locking overhead. The backend infrastructure, running on a Kubernetes cluster, ingests sub-second container telemetry via Prometheus and cAdvisor. Utilizing this telemetry, the system applies Particle Swarm Optimization (PSO) for in-place vertical rightsizing of container CPU quotas and memory boundaries without restarting pods. Simultaneously, Ant Colony Optimization (ACO) balances pod allocations across worker nodes, Dynamic Voltage and Frequency Scaling (DVFS) modulates processor frequencies during developer idle periods, and an Adaptive Resource \& Knowledge Lifecycle Manager (ARKLM) snapshots dormant workspaces.

Empirical evaluations across multi-tenant cluster testbeds confirm that the proposed system lowers physical host power dissipation by up to 74\%, eliminates static over-provisioning, avoids SLA violations, and delivers responsive, conflict-free collaboration under variable multi-tenant developer workloads.

\vspace{0.4cm}
\textbf{Keywords:} Cloud IDE, Real-Time Collaboration, CRDT, Kubernetes, Particle Swarm Optimization, Ant Colony Optimization, DVFS, Green Computing, Vertical Autoscaling.

% =========================================================
% ACKNOWLEDGMENT
% =========================================================
\clearpage
\chapter*{Acknowledgment}
\thispagestyle{plain}
\addcontentsline{toc}{front}{Acknowledgment}

\noindent
It gives us immense pleasure and satisfaction in presenting this preliminary project report on \textbf{`Resource-Aware Collaborative Cloud IDE'}.

\vspace{0.4cm}
We express our deepest and most sincere gratitude to our internal guide, \textbf{Prof. Yogita Hambir}, for her continuous intellectual direction, regular reviews, and constructive critiques throughout the inception, design, and experimental evaluation of this work. Her guidance provided the clarity necessary to establish our experimental methodologies.

\vspace{0.4cm}
We convey our heartfelt appreciation to our Project Coordinator, \textbf{Prof. Anup Kadam}, and to \textbf{Prof. (Dr.) Sunil Dhore}, Head of the Department of Computer Engineering, for facilitating essential server hardware access, container testbeds, and an encouraging academic environment. We also express our gratitude to our Principal, \textbf{Dr. B. P. Patil}, for supporting project research at the Army Institute of Technology, Pune.

\vspace{0.4cm}
Finally, we extend our sincere thanks to all our departmental faculty members, technical laboratory assistants, and peers whose encouragement and support contributed directly to the completion of this report.

\vfill
\begin{flushright}
\begin{tabular}{ll}
    \textbf{Abhay Joshi} & (Exam No: 7402) \\
    \textbf{Anish Kumar Kaler} & (Exam No: 7411) \\
    \textbf{Anish Kumar Pandey} & (Exam No: 7412) \\
    \textbf{Bikash Sharma} & (Exam No: 7418) \\[0.2cm]
    \multicolumn{2}{r}{\small B.E. (Computer Engineering)} \\
    \multicolumn{2}{r}{\small Department of Computer Engineering} \\
    \multicolumn{2}{r}{\small Army Institute of Technology, Pune}
\end{tabular}
\end{flushright}
\vspace{0.5cm}

\clearpage
\tableofcontents

\clearpage
\addcontentsline{toc}{front}{List of Figures}
\listoffigures

\clearpage
\addcontentsline{toc}{front}{List of Tables}
\listoftables

\clearpage
\pagenumbering{arabic}
\setcounter{page}{1}

% =========================================================
% CHAPTER 1: INTRODUCTION
% =========================================================
\chapter{INTRODUCTION}

\section{Details of Project Work}
Software engineering practices have transitioned steadily from local desktop development environments toward cloud-native web architectures. Historically, setting up a development workspace required developers to manually install and configure operating system-specific compilers, language runtimes, local database servers, native system libraries, and custom debugger tools. This localized approach created significant operational friction, leading to environment inconsistencies---frequently referred to as the ``works on my machine'' dilemma---where code working locally encounters bugs or failures in production due to minor differences in operating system builds or library versions.

Cloud-based Integrated Development Environments (Cloud IDEs) address this problem by moving the developer's execution sandbox into cloud containers and providing an editor interface via web browsers. Developers can access pre-configured workspaces using a browser link, eliminating setup overhead. However, operating multi-tenant cloud IDE infrastructure introduces significant engineering challenges:
\begin{enumerate}
    \item \textbf{High Synchronization Overhead:} Concurrent collaborative document editing has traditionally relied on Operational Transformation (OT). OT requires central server mediation to serialize operations and maintain total order, resulting in communication bottlenecks and noticeable keystroke latency in geographically distributed teams.
    \item \textbf{Resource Over-Provisioning:} Developer workflows are bursty. Developers spend considerable time reading documentation, designing code, or typing, which requires negligible CPU and memory (often below 5\% utilization). In contrast, compiling code, resolving dependencies, and running test suites create high, transient compute spikes. Conventional cloud IDEs address this by statically sizing container limits to peak demands. Consequently, compute and memory allocations remain largely idle during typing and reading, driving up cloud infrastructure costs and energy consumption.
\end{enumerate}

This project presents the design and implementation of a \textbf{Resource-Aware Collaborative Cloud IDE} that links frontend editing activity with backend infrastructure scaling. The editor uses decentralized Conflict-free Replicated Data Types (CRDTs), ensuring that concurrent peer edits converge deterministically with sub-100ms peer synchronization latency without central locks. Simultaneously, the underlying Kubernetes cluster operates an active telemetry and optimization pipeline: Particle Swarm Optimization (PSO) calculates in-place vertical CPU and memory quotas; Ant Colony Optimization (ACO) balances pod allocations across nodes to avoid thermal hotspots; Dynamic Voltage and Frequency Scaling (DVFS) throttles physical processor frequencies during idle editing intervals; and an Adaptive Resource \& Knowledge Lifecycle Manager (ARKLM) snapshots inactive pods to persistent storage.

\section{Objectives}
The primary objective of this project is to optimize computing resource utilization and reduce energy dissipation in a multi-tenant cloud development platform while maintaining sub-100ms interactive collaboration latencies. The specific engineering objectives are:
\begin{enumerate}
    \item \textbf{Decentralized Real-Time Collaborative Synchronization:} Implement an operation-based sequence CRDT architecture using Yjs and WebSockets, ensuring eventual consistency and deterministic conflict resolution without centralized operational locks.
    \item \textbf{Lightweight Sandbox Provisioning:} Develop an automated container orchestration engine on Kubernetes capable of launching isolated developer execution pods within 5 seconds of request.
    \item \textbf{Sub-Second Telemetry Ingestion:} Build an observability pipeline that scrapes container-level CPU, memory, and I/O metrics at 1-second intervals via cAdvisor and Prometheus.
    \item \textbf{Vertical Container Rightsizing via PSO:} Implement a continuous Particle Swarm Optimization algorithm that dynamically rightsizes container CPU and RAM limits to active demand without requiring container restarts.
    \item \textbf{Heuristic Cluster Balancing via ACO:} Implement an Ant Colony Optimization scheduler to assign workspace pods across cluster nodes, balancing load and preventing thermal hotspots.
    \item \textbf{Hardware-Level Power Modulation via DVFS:} Implement a host daemon that modulates physical CPU core frequencies based on developer interaction states, lowering idle power consumption while meeting SLA commitments.
    \item \textbf{Automated Workspace Lifecycle Management:} Implement ARKLM to track developer think-time, transition idle pods into low-frequency states, and snapshot inactive workspaces to persistent storage after 15 minutes of inactivity.
\end{enumerate}

\section{Scope of the Project}
The defined scope of this project includes:
\begin{itemize}
    \item \textbf{Web Client Interface:} A single-page application built on React 18, featuring the Monaco Editor engine, an xterm.js terminal emulator, and Yjs CRDT synchronization.
    \item \textbf{Platform Control Plane:} Microservices managing user authentication, room orchestration, Git repository ingestion, and WebSocket message routing.
    \item \textbf{Optimization Engine:} Metaheuristic PSO and ACO scheduling controllers integrated with host-level Linux cpufreq DVFS drivers.
    \item \textbf{Target Infrastructure:} A multi-node Kubernetes cluster utilizing Linux cgroups v2 and containerd.
\end{itemize}

The project scope excludes native desktop applications, GPU hardware acceleration for deep learning workloads, multi-cloud cross-provider federation, and enterprise billing/metering subsystems.

\section{Motivation of the Project}
Three core operational insights motivate this research:
\begin{itemize}
    \item \textbf{Workload Burstiness in Software Engineering:} Developer profiling indicates that active compilation and test runs account for less than 8\% of a typical interactive programming session. The remaining 92\% consists of code comprehension, reading, and text entry. Static resource provisioning sized for peak compilation demands wastes substantial server capacity.
    \item \textbf{Semantic Awareness Gap in Standard Autoscalers:} Generic Kubernetes autoscalers (such as HPA and VPA) lack insight into application-level user context. They cannot distinguish between an idle terminal and an active developer thinking between keystrokes. Coupling IDE interaction telemetry directly with infrastructure rightsizing enables proactive and fine-grained resource modulation.
    \item \textbf{Environmental and Economic Imperatives of Green Computing:} Data centers worldwide consume substantial electrical energy. By dynamically reducing operating frequencies and reclaiming dormant memory allocations, cloud IDE providers can lower operational costs and carbon emissions.
\end{itemize}

% =========================================================
% CHAPTER 2: LITERATURE SURVEY
% =========================================================
\chapter{LITERATURE SURVEY}

Managing multi-tenant cloud development environments requires addressing two distinct engineering challenges: infrastructure resource optimization and low-latency distributed collaboration. This chapter surveys foundational and recent works across both domains.

\section{Cluster Management and Container Scaling Systems}
Verma et al.~\cite{borg} introduce Borg, Google's production cluster management architecture. Borg pioneered aggressive resource over-commitment, bin-packing, and priority-based task preemption. While Borg demonstrated that dynamic task packing improves utilization, it treats tasks as opaque black boxes. It lacks visibility into fine-grained interactive developer states and does not interface with CPU frequency scaling or collaborative synchronization fabrics.

Shan et al.~\cite{aras} propose the Adaptive Resource Allocation System (ARAS) for containerized workflows on Kubernetes. ARAS implements a MAPE-K loop for vertical container scaling. However, ARAS relies on reactive threshold rules, introducing scaling lag during sharp transitions.

Al-Dhuraibi et al.~\cite{pace} develop the PACE framework, which pairs Convolutional Neural Networks (CNN) with K-means clustering for predictive auto-scaling. Evaluated against Redis NoSQL containers under the YCSB benchmark, PACE accurately forecasted request spikes and dynamically scaled virtual CPU allocations. However, PACE was restricted to single-tier NoSQL workloads and did not evaluate multi-resource vertical rightsizing or interactive think-time within developer environments.

Toka et al.~\cite{sata} present the SATA algorithm, an SLA-adaptive moving-average threshold controller for the Kubernetes Horizontal Pod Autoscaler. SATA dynamically adjusts CPU scaling thresholds under bursty access logs (such as WorldCup98 and SEC EDGAR traces) to prevent SLO violations without complex deep learning models.

\section{Metaheuristic and Evolutionary Scheduling in Cloud Environments}
To address the NP-hard nature of multi-objective container scheduling, researchers have investigated metaheuristic algorithms. Dubey and Sharma~\cite{psoaco} formulate a hybrid PSO-ACO algorithm for cloud task scheduling. Evaluated in CloudSim, their approach achieved faster makespan convergence than standalone PSO or ACO. However, their work was confined to static simulations, lacking validation against live containerized runtimes or live telemetry.

Hsu et al.~\cite{pso_lstm} present KIF, integrating Particle Swarm Optimization with Long Short-Term Memory (LSTM) networks for predictive pod scheduling on edge-cloud nodes, achieving an RMSE of $\approx 0.28$ on ENACT telemetry traces. Rajesh et al.~\cite{cshbo} develop CS-HBO, combining Cuckoo Search with Elephant Herd and Spider Monkey Optimization to minimize network distance and system failure rates across physical machines. Similarly, Al-Janabi and Mahmood~\cite{atsoa} introduce ATSOA, an adaptive tree seed optimization algorithm with genetic operators that reduces execution costs for scientific workflows on heavy virtual machines.

Wang et al.~\cite{tspra} present TSPRA, combining tabular diffusion data augmentation (TabDDPM) with Mamba sequence modeling and an ACO hyper-heuristic over Whale Optimization and Quantum PSO. Evaluated across DeathStarBench microservices, TSPRA eliminated resource over-provisioning while maintaining QoS, but did not incorporate interactive developer idle states or hardware-level clock modulation. In related work, Conlloovia~\cite{conlloovia} implements Mixed Integer Linear Programming (MILP) to find opportunistic, cost-optimal resource allocations for elastic cloud applications. ROCAMS~\cite{rocams} evaluates container consolidation using adapted genetic algorithms and PSO, minimizing replacement count by up to 130\% while conserving power. Furthermore, DAScheduler~\cite{dascheduler} demonstrates dependency-aware scheduling to achieve load balancing across multi-zone container topologies.

\section{Energy Optimization and Hardware Frequency Scaling}
Energy efficiency is a primary design goal in modern cloud infrastructure. Beena et al.~\cite{dvfs} combine ACO, PSO, and Dynamic Voltage and Frequency Scaling (DVFS), demonstrating host power reductions of up to 28\% during low-compute periods. However, their workload model was oriented around batch processing jobs and did not use interactive developer activity to trigger frequency changes.

Krasteva et al.~\cite{petas} propose PETAS, integrating Random Forest task-duration predictions with Kubernetes autoscaling, reducing cumulative pod uptime by up to 50\%. Li et al.~\cite{cec} formulate CEC, utilizing Fast Fourier Transform (FFT) pattern recognition and self-adaptive PID controllers for edge container provisioning, achieving 87.52\% prediction accuracy on Alibaba traces. LBSM~\cite{lbsm} optimizes container placement on OpenStack by evaluating intranode and internode imbalance under NUMA constraints. CARSS~\cite{carss} applies formal model checking to prove context-aware cloud resource scheduling bounds, while Qafhe~\cite{qafhe} provides dynamic benchmarking of Kubernetes orchestration performance.

\section{Collaborative Editing and Synchronization Algorithms}
Real-time collaborative editing has evolved from centralized locking mechanisms toward decentralized consensus. Sun et al.~\cite{crdtsurvey} survey the transition from Operational Transformation (OT) to Conflict-free Replicated Data Types (CRDTs). While OT requires centralized servers to order operations and complex transformation matrices, CRDTs achieve mathematically verified eventual consistency through commutative operations. Integrating CRDT synchronization state with cloud infrastructure autoscalers forms a key contribution of this project.

\clearpage
\section{Literature Survey Summary Table}
Table~\ref{tab:litsum} provides a comparative summary of the surveyed research works.

\begin{table}[htbp]
\centering
\scriptsize
\caption{Comprehensive Literature Survey Summary Table}
\label{tab:litsum}
\begin{tabularx}{\linewidth}{@{}p{0.8cm}p{2.8cm}p{2.3cm}p{3.2cm}X@{}}
\toprule
\textbf{Ref.} & \textbf{Method / Technique} & \textbf{Dataset / Trace} & \textbf{Evaluation Metrics} & \textbf{Targeted Research Gap} \\ \midrule
\cite{borg} & Over-commitment \& task packing & Google Production Logs & Cluster utilization, preemption rate & Workloads treated as black boxes; no interactive state awareness. \\
\cite{aras} & ARAS: MAPE-K loop on K8s & Live K8s Pods & CPU/RAM limits, response time & Reactive rules cause scaling lag; lacks metaheuristic search. \\
\cite{pace} & PACE: CNN + K-means clustering & GCP Redis, YCSB & CPU \%, vCPU alloc, throughput & Single-tier NoSQL focus; lacks multi-resource vertical rightsizing. \\
\cite{sata} & SATA: Moving-average HPA & WorldCup98, EDGAR logs & SLO compliance \%, pods count & Purely reactive horizontal scaling; requires violation window. \\
\cite{psoaco} & Hybrid PSO-ACO task scheduling & CloudSim simulation & Makespan, convergence speed & Static simulation only; not validated on live containers. \\
\cite{pso_lstm} & KIF: PSO + LSTM pod scheduling & ENACT edge dataset & RMSE, MAE, MAPE & Offline telemetry; lacks live in-place vertical rightsizing. \\
\cite{cshbo} & CS-HBO: Cuckoo Search + Elephant & Simulated PM clusters & Total Distance, Cost, SF & Synthetic static topologies; unvalidated on live clusters. \\
\cite{atsoa} & ATSOA: Tree Seed Algorithm & Montage, CyberShake & Overall Cost, Makespan, Util & Targeted at heavy VMs rather than lightweight containers. \\
\cite{tspra} & TSPRA: TabDDPM + Mamba & DeathStarBench & Latency, MSE, Utilization & Evaluated on microservice DAGs; lacks idle think-time model. \\
\cite{petas} & PETAS: Random Forest auto-scaling & Optimization benchmarks & Pod uptime, optimization time & Requires static initial generation; lacks think-time scaling. \\
\cite{cec} & CEC: FFT pattern + PID control & Alibaba Cluster Traces & Prediction accuracy, latency & Tailored for periodic loads; unsuited for stochastic bursts. \\
\cite{dvfs} & ACO + PSO with kernel DVFS & Batch server tasks & Host power (W), energy saving \% & Generic batch workloads; lacks real-time think-time triggers. \\
\cite{conlloovia} & Conlloovia: Opportunistic MILP & Alibaba ML job traces & Monetary cost, QoS violation \% & High solver compute times; assumes static scheduling windows. \\
\cite{rocams} & ROCAMS: EcoBalancer GA/PSO & Simulated IaaS data center & Power savings, replacement count & Simulated VM consolidation; lacks network mesh routing. \\
\cite{dascheduler} & DAScheduler: Dependency-aware & Synthetic cloud topology & Load balance, network traffic & Synthetic dependency matrix; lacks fluctuating interactive load. \\
\cite{lbsm} & LBSM: NUMA load-balancer & OpenStack C1--C9 specs & Internode/intranode imbalance & Limited to OpenStack IaaS; no native K8s pod support. \\
\cite{crdtsurvey} & Algorithmic survey: OT vs CRDTs & Theoretical models & Convergence, complexity & Purely algorithmic; decoupled from cloud orchestration. \\ \bottomrule
\end{tabularx}
\end{table}

% =========================================================
% CHAPTER 3: REQUIREMENT ANALYSIS
% =========================================================
\chapter{REQUIREMENT ANALYSIS}

\section{Software Requirements Specification (SRS)}
\subsection{Functional Requirements}
\begin{itemize}
    \item \textbf{FR1 - Automated Workspace Provisioning:} The platform shall dynamically instantiate, configure, and launch an isolated containerized workspace pod within 5 seconds of a user request.
    \item \textbf{FR2 - Real-Time Collaborative Synchronization:} The system shall synchronize concurrent document insertions and deletions across all connected peers with a propagation latency under 100ms, using operation-based sequence CRDTs.
    \item \textbf{FR3 - Sub-Second Telemetry Ingestion:} The telemetry daemon shall collect container-level CPU usage, memory resident set size (RSS), and disk I/O metrics at 1-second intervals, maintaining a 60-second sliding evaluation window.
    \item \textbf{FR4 - Vertical Container Rightsizing:} The optimization core shall continuously evaluate container telemetry and adjust container CPU core quotas and memory limits in place via PSO, without restarting the pod.
    \item \textbf{FR5 - Dynamic Pod Placement:} The cluster scheduler shall run an ACO algorithm to assign incoming workspace pods across cluster nodes, minimizing utilization variance and preventing hotspots.
    \item \textbf{FR6 - Hardware-Level Energy Scaling (DVFS):} The host daemon shall modulate physical CPU frequencies between $f_{\text{low}}$ and $f_{\max}$ based on active developer interactions, lowering idle power consumption.
    \item \textbf{FR7 - Workspace Snapshotting and Eviction (ARKLM):} Workspaces remaining idle (no keystrokes or terminal I/O) for more than 15 minutes shall be snapshotted to persistent storage and their compute resources evicted.
    \item \textbf{FR8 - Presence and Cursor Tracking:} The collaborative engine shall broadcast cursor positions, text selections, and user presence markers to all active participants in a room in real time.
\end{itemize}

\subsection{Non-Functional Requirements}
\begin{itemize}
    \item \textbf{NFR1 - Performance and Scalability:} Peer edit propagation latency must remain below 100ms for up to 10 concurrent collaborators per file buffer, with the client editor maintaining 60 frames per second during typing.
    \item \textbf{NFR2 - Security and Isolation:} Container sandboxes must execute as non-root users, bounded by Linux cgroups v2, distinct namespaces, and strict seccomp system call filters.
    \item \textbf{NFR3 - Availability and Fault Tolerance:} The edge gateway and WebSocket broker shall maintain 99.5\% service uptime. If a network disruption occurs, the client CRDT engine must support offline editing and resynchronize cleanly upon reconnection.
    \item \textbf{NFR4 - Maintainability and Extensibility:} The backend must adopt a modular microservices architecture, with communication handled via standard REST and WebSocket protocols.
\end{itemize}

% =========================================================
% CHAPTER 4: MATHEMATICAL MODELING
% =========================================================
\chapter{ALGORITHM ANALYSIS AND MATHEMATICAL MODELING}

\section{Collaborative Synchronization Algorithm (CRDT)}
To avoid centralized locking, the collaborative editing layer uses an operation-based sequence Conflict-free Replicated Data Type (CRDT). Let document state $S$ be represented as an ordered sequence of unique atomic character structures:
\begin{equation}
    S = \langle a_1, a_2, \dots, a_n \rangle
\end{equation}
Each atom $a_i$ is uniquely identified by a tuple:
\begin{equation}
    a_i = \langle \text{client\_id}, \text{clock\_tick}, \text{char\_val}, \text{origin\_left}, \text{origin\_right} \rangle
\end{equation}
where $\text{clock\_tick}$ represents a monotonically increasing Lamport logical clock, $\text{client\_id}$ breaks ties deterministically, and $(\text{origin\_left}, \text{origin\_right})$ specify the spatial insertion bounds. For any two concurrent operations $o_1$ and $o_2$, strong eventual consistency is guaranteed by commutative operation ordering:
\begin{equation}
    S \circ o_1 \circ o_2 = S \circ o_2 \circ o_1
\end{equation}
Character deletions are implemented using tombstones rather than physical deletions. This preserves the structural coordinate system for concurrent operations without centralized arbitration.

\section{Particle Swarm Optimization (PSO) for Vertical Rightsizing}
The vertical scaling controller formulates container resource allocation as a multi-objective cost minimization problem. Let $\mathbf{X}_i = [c_i^{\text{alloc}}, m_i^{\text{alloc}}]$ denote the continuous allocation vector for container $i$, representing allocated CPU cores and memory limits. The cost objective function is formulated as:
\begin{equation}
    F_{\text{cost}}(\mathbf{X}) = \sum_{i=1}^N \left( w_1 \frac{c_i^{\text{alloc}} - c_i^{\text{used}}}{c_i^{\text{alloc}}} + w_2 \frac{m_i^{\text{alloc}} - m_i^{\text{used}}}{m_i^{\text{alloc}}} \right) + \lambda P_{\text{SLA}}
\end{equation}
where $w_1$ and $w_2$ are weighting parameters balancing CPU and memory slack, and $P_{\text{SLA}}$ is a penalty function activated when utilization exceeds safe operating thresholds $\theta_{\text{cpu}}$ and $\theta_{\text{mem}}$:
\begin{equation}
    P_{\text{SLA}} = \sum_{i=1}^N \max\left(0, \frac{c_i^{\text{used}} - \theta_{\text{cpu}} c_i^{\text{alloc}}}{c_i^{\text{alloc}}}\right)^2 + \max\left(0, \frac{m_i^{\text{used}} - \theta_{\text{mem}} m_i^{\text{alloc}}}{m_i^{\text{alloc}}}\right)^2
\end{equation}
Each particle $k$ in the swarm tracks its position $\mathbf{x}_k$ and velocity $\mathbf{v}_k$ across iterations $t$. Kinematic updates follow:
\begin{equation}
    \mathbf{v}_k^{t+1} = \omega \mathbf{v}_k^t + c_1 r_1 (\mathbf{p}_{\text{best},k}^t - \mathbf{x}_k^t) + c_2 r_2 (\mathbf{g}_{\text{best}}^t - \mathbf{x}_k^t)
\end{equation}
\begin{equation}
    \mathbf{x}_k^{t+1} = \mathbf{x}_k^t + \mathbf{v}_k^{t+1}
\end{equation}
Here, $\omega$ is the dynamic inertia weight, $c_1$ and $c_2$ denote cognitive and social acceleration coefficients, and $r_1, r_2 \sim \mathcal{U}(0,1)$. The global optimal vector $\mathbf{g}_{\text{best}}$ is mapped directly to Kubernetes cgroups v2 control parameters (\texttt{cpu.max} and \texttt{memory.max}).

\section{Ant Colony Optimization (ACO) for Pod Placement}
Pod placement across heterogeneous worker nodes is handled via Ant Colony Optimization. The probability $P_{ij}$ of scheduling pod $i$ on worker node $j$ is governed by:
\begin{equation}
    P_{ij} = \frac{[\tau_{ij}]^\alpha [\eta_{ij}]^\beta}{\sum_{k \in \mathcal{A}} [\tau_{ik}]^\alpha [\eta_{ik}]^\beta}
\end{equation}
where $\tau_{ij}$ represents the historical pheromone concentration, and $\eta_{ij}$ is the load balance heuristic:
\begin{equation}
    \eta_{ij} = \frac{1}{1 + |U_j^{\text{CPU}} - \bar{U}^{\text{CPU}}| + |U_j^{\text{MEM}} - \bar{U}^{\text{MEM}}|}
\end{equation}
Here, $\bar{U}^{\text{CPU}}$ and $\bar{U}^{\text{MEM}}$ represent the cluster mean resource utilizations. Pheromone evaporation and deposition follow:
\begin{equation}
    \tau_{ij} \leftarrow (1 - \rho)\,\tau_{ij} + \sum_{k=1}^A \Delta \tau_{ij}^k, \quad \Delta \tau_{ij}^k = \frac{Q}{1 + \text{Var}(U^k)}
\end{equation}

\section{Dynamic Voltage and Frequency Scaling (DVFS)}
Digital CMOS processor power dissipation is governed by static and dynamic components:
\begin{equation}
    P_{\text{total}} = P_{\text{static}} + P_{\text{dynamic}} = I_{\text{leakage}} V + \alpha C_L V^2 f
\end{equation}
Because operational supply voltage $V$ scales approximately linearly with operating frequency $f$ within stable margins, dynamic power scales cubically:
\begin{equation}
    P_{\text{dynamic}} \propto f^3
\end{equation}
When the telemetry agent detects low CPU demand ($U_i^{\text{CPU}} < \theta_{\text{idle}}$) over a sustained evaluation window, the CPU core frequency is shifted to $f_{\text{low}}$ via the Linux \texttt{userspace} governor. Frequencies are restored to $f_{\max}$ upon user keystrokes or compilation execution requests.

% =========================================================
% CHAPTER 5: DETAILED DESIGN
% =========================================================
\chapter{DETAILED DESIGN}

\section{Architectural Design}
The system architecture follows a layered microservices pattern designed to decouple collaborative text synchronization from cluster orchestration and hardware optimization. The architecture organizes communication flow into distinct, loosely coupled layers:

\subsection{Client and Edge Interaction Layer}
The user interface executes as a browser-native Single Page Application (SPA) built using React 18. It integrates the Monaco Editor engine to provide IDE features such as syntax highlighting, language server protocol (LSP) diagnostics, code completion, and multi-file project tree management. For interactive terminal sessions, the client incorporates an xterm.js terminal emulator, communicating with container processes over secure WebSockets (WSS).

All client traffic terminates at an NGINX API Gateway that manages TLS certificates, enforces rate limits, handles CORS policies, and validates JSON Web Tokens (JWT). The gateway inspects incoming HTTP connection headers and transparently upgrades collaborative editing and terminal traffic to full-duplex WebSocket connections.

\subsection{State Engine and Collaboration Fabric}
Collaborative state synchronization bypasses traditional database transactions. Instead, document modifications are captured as compact binary deltas by the client-side Yjs CRDT library. These deltas are transmitted over WebSockets to Node.js \texttt{y-websocket} workers.

To scale across multiple collaboration workers, instances are connected through a Redis 7.2 Pub/Sub message bus. When a developer modifies code in a room, the local worker applies the commutative delta to its in-memory document state and publishes the update to a room-specific channel (\texttt{room:<roomId>:events}). All worker nodes hosting active connections for that room receive the delta and broadcast it to connected clients. Cursor tracking and presence awareness are managed using transient Redis hash sets, providing peer visibility without adding write overhead to persistent storage.

\subsection{High-Frequency Telemetry Pipeline}
The optimization layer relies on accurate, high-frequency metrics. Standard Kubernetes metric pipelines scrape data at 15- to 60-second intervals, which is too coarse to capture short developer compilation bursts. To address this, cAdvisor daemons run on each worker node to sample container metrics at 1-second intervals.

A local Prometheus instance scrapes these endpoints and maintains a 60-second sliding evaluation window. This sliding window provides the optimization engine with CPU core utilization, memory Resident Set Size (RSS), page cache metrics, and disk I/O rates.

\subsection{Swarm Optimization and Power Governance}
The optimization core runs as a control daemon on the Kubernetes master node. At each evaluation step, it retrieves the sliding-window metrics and executes two metaheuristic algorithms:
\begin{enumerate}
    \item \textbf{Vertical Rightsizing (PSO):} The PSO engine evaluates the multi-objective fitness function across the current container population. If slack or saturation exceeds predefined safety thresholds, the swarm computes an updated allocation vector $\mathbf{g}_{\text{best}}$. The engine then applies these limits directly to the Linux cgroups v2 controllers via the Kubernetes pod patch API without restarting containers.
    \item \textbf{Cluster Balancing (ACO):} When new workspaces are initialized, the ACO scheduler calculates a balanced placement distribution based on real-time node metrics and pheromone weights, avoiding resource-constrained nodes.
\end{enumerate}

Simultaneously, a host-level DVFS daemon monitors node-level utilization. If all containers on a worker node remain idle beyond the think-time threshold, the daemon shifts physical CPU core frequencies to $f_{\text{low}}$ using the Linux \texttt{cpufreq} interface, scaling back up when new activity is detected.

\clearpage
\subsection{System Architecture Diagram}

\begin{figure}[H]
\centering
\singlespacing
\pngdiag{arch_diag.png}{%
\begin{tikzpicture}[font=\footnotesize, xscale=1.05, yscale=1.05,
  dn/.style={-{Stealth[length=2.6mm,width=2.4mm]}, very thick, black}]

\node[rounded corners=4pt,draw=black,thick,fill=white,minimum width=11cm,minimum height=1.0cm,
      align=center,font=\normalsize\bfseries] (ttl) at (0,-0.5)
  {Proposed System Architecture\\[-1pt]{\small\mdseries (Resource-Aware Collaborative Cloud IDE)}};

\archlayer{black}{-1.5}{-3.7}{Client Layer (Browser)}
\archlayer{black}{-4.2}{-6.4}{Edge and Platform Services Layer}
\archlayer{black}{-6.9}{-9.1}{Collaboration and State Engine Layer}
\archlayer{black}{-9.6}{-11.8}{Monitoring and Telemetry Layer}
\archlayer{black}{-12.3}{-16.9}{Resource Optimization Layer (Closed-Loop Controller)}
\archlayer{black}{-17.4}{-19.6}{Kubernetes Workspace Cluster}
\archlayer{black}{-20.1}{-22.5}{Environment Feedback (Performance Metrics)}

\archbox{black}{-4.2}{-2.85}{3.8}{1.3}{\textbf{Editor}\\React 18 SPA,\\Monaco Editor}
\archbox{black}{0}{-2.85}{3.8}{1.3}{\textbf{Terminal}\\xterm.js over WSS,\\stdout / stderr stream}
\archbox{black}{4.2}{-2.85}{3.8}{1.3}{\textbf{CRDT Client}\\Yjs, cursors and\\presence}

\archbox{black}{-4.65}{-5.55}{2.9}{1.3}{\textbf{NGINX Ingress}\\TLS, WSS upgrade,\\rate limiting}
\archbox{black}{-1.55}{-5.55}{2.9}{1.3}{\textbf{Auth Service}\\OAuth2 / JWT}
\archbox{black}{1.55}{-5.55}{2.9}{1.3}{\textbf{Workspace API}\\session provisioning}
\archbox{black}{4.65}{-5.55}{2.9}{1.3}{\textbf{Repo Service}\\clone, stack detection}

\archbox{black}{-4.65}{-8.25}{2.9}{1.3}{\textbf{y-websocket}\\CRDT delta sync}
\archbox{black}{-1.55}{-8.25}{2.9}{1.3}{\textbf{Redis Pub/Sub}\\room broadcast}
\archbox{black}{1.55}{-8.25}{2.9}{1.3}{\textbf{Presence Store}\\cursors, users}
\archbox{black}{4.65}{-8.25}{2.9}{1.3}{\textbf{PostgreSQL}\\users, SLA logs}

\archbox{black}{-4.2}{-10.95}{3.8}{1.3}{\textbf{cAdvisor}\\container CPU, RAM, I/O}
\archbox{black}{0}{-10.95}{3.8}{1.3}{\textbf{Prometheus}\\1 s scrape,\\60 s sliding window}
\archbox{black}{4.2}{-10.95}{3.8}{1.3}{\textbf{Node Metrics}\\queue length, power (W)}

\archbox{black}{0}{-13.35}{11.6}{0.8}{\textbf{State:} workload demand, current allocation, utilization, SLA status, idle time}
\draw[dn] (0,-13.75) -- (0,-14.2);
\archbox{black}{-4.2}{-14.85}{3.8}{1.3}{\textbf{PSO Rightsizing}\\vertical CPU / RAM\\under SLA limits}
\archbox{black}{0}{-14.85}{3.8}{1.3}{\textbf{ACO Scheduler}\\balanced pod-to-node\\placement}
\archbox{black}{4.2}{-14.85}{3.8}{1.3}{\textbf{DVFS + ARKLM}\\frequency scaling,\\idle reclamation}
\draw[dn] (0,-15.5) -- (0,-15.8);
\archbox{black}{0}{-16.2}{11.6}{0.8}{\textbf{Action:} resize cgroup limits, place / migrate pods, set CPU frequency, snapshot and reclaim idle workspaces}

\archbox{black}{-4.2}{-18.75}{3.8}{1.3}{\textbf{Worker Node 1}\\Workspace Pods\\(Node.js 20 / Python 3.11)}
\archbox{black}{0}{-18.75}{3.8}{1.3}{\textbf{Worker Node 2}\\Workspace Pods\\cgroups v2, containerd}
\archbox{black}{4.2}{-18.75}{3.8}{1.3}{\textbf{Worker Node 3}\\Kubelet + DVFS agent,\\persistent volumes}

\archbox{black}{-4.2}{-21.0}{3.8}{0.6}{Actual CPU utilization}
\archbox{black}{0}{-21.0}{3.8}{0.6}{Actual memory utilization}
\archbox{black}{4.2}{-21.0}{3.8}{0.6}{Sync latency (ms)}
\archbox{black}{-4.2}{-21.8}{3.8}{0.6}{SLA violation (\%)}
\archbox{black}{0}{-21.8}{3.8}{0.6}{Host power (W)}
\archbox{black}{4.2}{-21.8}{3.8}{0.6}{Cost and resource wastage}

\node[rounded corners=3pt,draw=black,thick,fill=white,minimum width=7.4cm,minimum height=0.9cm,
      align=center,font=\scriptsize] (rew) at (0,-23.6)
  {\textbf{Reward} (based on $F_{\text{cost}}$: resource waste, power, SLA penalty)};

\draw[dn] (0,-1.0)   -- (0,-1.5);
\draw[dn] (0,-3.7)  -- (0,-4.2);
\draw[dn] (0,-6.4)  -- (0,-6.9);
\draw[dn] (0,-9.1)  -- (0,-9.6);
\draw[dn] (0,-11.8) -- (0,-12.3);
\draw[dn] (0,-16.9) -- (0,-17.4);
\draw[dn] (0,-19.6) -- (0,-20.1);
\draw[dn] (0,-22.5) -- (rew.north);

\draw[dn] (rew.east) -| (7.6,-14.85) -- (6.6,-14.85);
\node[rotate=90,font=\tiny,fill=white,inner sep=1.5pt] at (7.6,-19.2) {feedback loop: metrics and reward};
\end{tikzpicture}}
\caption{Layered System Architecture of the Resource-Aware Collaborative Cloud IDE}
\label{fig:arch}
\end{figure}

\clearpage
\section{UML Design Diagrams}

\subsection{Use Case Diagram}
\begin{figure}[H]
\centering
\singlespacing
\pngdiag{usecase_diag.png}{%
\begin{tikzpicture}[xscale=1.05, yscale=1.05,
  uc/.style={ellipse, draw=black, thick, fill=white, text=black, text width=2.6cm, align=center,
             minimum height=1.0cm, inner sep=2pt, font=\footnotesize},
  inc/.style={-{Straight Barb[length=2.4mm,width=2.6mm]}, dashed, thick, black},
  rl/.style={font=\tiny\itshape, inner sep=1.5pt}]

\draw[thick] (2.3,1.4) rectangle (13.9,-19.3);
\node[font=\small\bfseries, fill=white, inner sep=3pt] at (8.1,1.4) {Resource-Aware Collaborative Cloud IDE};

\node[uc] (a1) at (5,0)      {Authenticate\\(Login / OAuth)};
\node[uc] (a2) at (5,-1.8)   {Create / Join Room};
\node[uc] (a3) at (5,-3.6)   {Clone Repository};
\node[uc] (a4) at (5,-5.4)   {Edit Code Collaboratively};
\node[uc] (a5) at (5,-7.2)   {View Presence \& Cursors};
\node[uc] (a6) at (5,-9.0)   {Run / Build Code};
\node[uc] (a7) at (5,-10.8)  {View Console Logs};
\node[uc] (a8) at (5,-14.0)  {Monitor Cluster Metrics};
\node[uc] (a9) at (5,-15.8)  {Configure Scaling Policies};
\node[uc] (a10) at (5,-17.6) {View Energy Reports};

\node[uc] (b1) at (11.2,-3.6)  {Analyze Dependencies\\Detect Stack};
\node[uc] (b2) at (11.2,-5.4)  {Resolve Conflicts\\(OT / CRDT)};
\node[uc] (b0) at (11.2,-7.2)  {Generate Dockerfile\\ImageSpec};
\node[uc] (b3) at (11.2,-9.0)  {Provision Container Pod};
\node[uc] (b4) at (11.2,-10.8) {Auto-Scale Pods\\(HPA / ARAA)};
\node[uc] (b5) at (11.2,-12.6) {Optimize Resources\\(PSO / ACO)};
\node[uc] (b6) at (11.2,-14.4) {Apply DVFS\\Frequency Scaling};
\node[uc] (b7) at (11.2,-16.2) {Reclaim Idle Sessions\\(ARKLM)};

\actor{dev}{0.2}{-5.4}{Developer}
\actor{adm}{0.2}{-15.8}{Cluster Admin}
\actor{git}{16.7}{-3.6}{\gl{external system}\\Git Host\\(GitHub / GitLab)}
\actor{kub}{16.7}{-12.6}{\gl{external system}\\Kubernetes\\Cluster API}

\foreach \n in {a1,a2,a3,a4,a5,a6,a7} { \draw[thick,black] (deve) -- (\n.west); }
\foreach \n in {a8,a9,a10} { \draw[thick,black] (adme) -- (\n.west); }
\draw[thick,black] (gitw) -- (b1.east);
\foreach \n in {b3,b4,b5,b6,b7} { \draw[thick,black] (kubw) -- (\n.east); }

\draw[inc] (a2.north) -- node[rl,right] {\gl{include}} (a1.south);
\draw[inc] (a3.east) -- node[rl,above] {\gl{include}} (b1.west);
\draw[inc] (a4.east) -- node[rl,above] {\gl{include}} (b2.west);
\draw[inc] (a6.east) -- node[rl,above] {\gl{include}} (b3.west);
\draw[inc] (b3.north) -- node[rl,right] {\gl{include}} (b0.south);
\draw[inc] (b4.north) -- node[rl,right] {\gl{extend}} (b3.south);
\draw[inc] (b4.south) -- node[rl,right] {\gl{include}} (b5.north);
\draw[inc] (b5.south) -- node[rl,right] {\gl{include}} (b6.north);
\end{tikzpicture}}
\caption{Use Case Diagram of the Resource-Aware Collaborative Cloud IDE}
\label{fig:usecase}
\end{figure}

Figure~\ref{fig:usecase} models the functional interactions between system actors and primary use cases. The Developer interacts with workspace-level features, including authentication, room creation, real-time editing, presence tracking, and terminal command execution. The Cluster Administrator accesses monitoring features such as telemetry dashboards, scaling thresholds, and energy reports. External systems (Git Hosts and Kubernetes) support backend operations, including repository cloning, pod provisioning, swarm-based scaling, DVFS frequency modulation, and session reclamation.

\clearpage
\subsection{Class Diagram: Collaborative Real-Time Workspace}
\begin{figure}[H]
\centering
\singlespacing
\pngonly{class_diag1.png}
\caption{Class Diagram: Real-Time Collaborative Engine}
\label{fig:class_collab}
\end{figure}

Figure~\ref{fig:class_collab} details the class hierarchy of the collaborative editing engine. The \texttt{CollabRoom} class manages connected user sessions and routes document operations. The \texttt{CRDTEngine} handles the atom sequence, Lamport clocks, and commutative insert/delete operations. The \texttt{PresenceTracker} records cursor line and column offsets to provide real-time peer visibility.

\clearpage
\subsection{Class Diagram: Workspace Provisioning and Runtime Execution}
\begin{figure}[H]
\centering
\singlespacing
\pngonly{class_diag2.png}
\caption{Class Diagram: Workspace Container Provisioning \& Runtime}
\label{fig:class_prov}
\end{figure}

Figure~\ref{fig:class_prov} models the classes responsible for workspace provisioning and container execution. The \texttt{WorkspaceService} coordinates with the \texttt{RepoAnalyzer} to infer language dependencies and generate an \texttt{ImageSpec}. The \texttt{PodManager} interfaces with \texttt{K8sClientWrapper} to create unprivileged container sandboxes, attach persistent storage volumes, and manage interactive terminal streams.

\clearpage
\subsection{Class Diagram: Heuristic Optimization and Lifecycle}
\begin{figure}[H]
\centering
\singlespacing
\pngonly{class_diag3.png}
\caption{Class Diagram: Heuristic Optimization \& Lifecycle Engine}
\label{fig:class_opt}
\end{figure}

Figure~\ref{fig:class_opt} specifies the classes comprising the telemetry and optimization pipeline. The \texttt{MetricsCollector} retrieves time-series data from Prometheus into rolling ring buffers. The \texttt{PSORightsizer} updates particle positions and velocities to optimize CPU/RAM allocations, while the \texttt{ACOScheduler} computes balanced pod placement. The \texttt{DVFSScalingController} and \texttt{ARKLMReclaimer} manage hardware frequency states and workspace snapshotting.

\clearpage
\subsection{Sequence Diagram: Session Setup and Real-Time Collaboration}
\begin{figure}[H]
\centering
\singlespacing
\pngdiagrot{seq_diag1.png}{%
\begin{tikzpicture}[xscale=1.1, yscale=1.2]
\def\xD{0} \def\xI{6.5} \def\xC{13} \def\xR{19.5}
\node[sobj] at (\xD,0) {\underline{Developer}\\{\tiny\gl{actor}}};
\node[sobj] at (\xI,0) {\underline{BrowserIDE}\\{\tiny(client)}};
\node[sobj] at (\xC,0) {\underline{CollabServer}\\{\tiny(WebSocket, OT / CRDT)}};
\node[sobj] at (\xR,0) {\underline{RepoEngine}};
\foreach \x in {\xD,\xI,\xC,\xR} { \draw[dashed,thick,black] (\x,-0.5) -- (\x,-13.5); }

\frag{-1.5}{-7.5}{14.6}{-11.6}{loop}{[session active]}{0.4}

\act{\xI}{-1.4}{-3.4} \act{\xI}{-4.4}{-6.8} \act{\xI}{-8.7}{-10.7}
\act{\xC}{-2.4}{-3.4} \act{\xC}{-9.7}{-10.7}
\act{\xR}{-4.4}{-6.8}

\sq{sm}{\xD}{\xI}{-1.4}{1: openRoom(roomId, repoUrl)}
\sq{sm}{\xI}{\xC}{-2.4}{2: wsConnect(roomId, jwt)}
\sq{sr}{\xC}{\xI}{-3.4}{3: ack + document snapshot}
\sql{sm}{\xI}{\xR}{-4.4}{4: cloneRepo(repoUrl)}{0.72}
\selfm{\xR}{-5.4}{5: inspectDependencies(),\\detect stack}{2.8}
\sql{sr}{\xR}{\xI}{-6.8}{6: ExecutionProfile (ImageSpec + Dockerfile)}{0.5}

\sq{sm}{\xD}{\xI}{-8.7}{7: type / edit code}
\sq{sm}{\xI}{\xC}{-9.7}{8: sendDelta(changeSet)}
\sq{sr}{\xC}{\xI}{-10.7}{9: broadcastSync(delta) to all peers}
\end{tikzpicture}}
\caption{Sequence Diagram, Part A: Room Join, Repository Analysis, and Collaborative Editing}
\label{fig:seq_a}
\end{figure}

Figure~\ref{fig:seq_a} documents the message flow during room initialization and collaborative editing. Operations include client authentication, WebSocket handshake, CRDT state hydration, automated Git repository cloning, runtime environment detection, and peer-to-peer edit broadcasting.

\clearpage
\subsection{Sequence Diagram: Execution, Scaling, and Teardown}
\begin{figure}[H]
\centering
\singlespacing
\pngdiag{seq_diag2.png}{%
\begin{tikzpicture}[xscale=1.0, yscale=1.1]
\def\xD{0} \def\xI{4} \def\xK{8} \def\xO{12} \def\xP{16}
\node[sobj] at (\xD,0) {\underline{Developer}\\{\tiny\gl{actor}}};
\node[sobj] at (\xI,0) {\underline{BrowserIDE}\\{\tiny(client)}};
\node[sobj] at (\xK,0) {\underline{K8sOrchestrator}};
\node[sobj] at (\xO,0) {\underline{Optimizer}\\{\tiny(PSO / ACO)}};
\node[sobj] at (\xP,0) {\underline{WorkerPod}\\{\tiny(DVFS-controlled)}};
\foreach \x in {\xD,\xI,\xK,\xO,\xP} { \draw[dashed,thick,black] (\x,-0.5) -- (\x,-22.5); }

\frag{6.5}{-7.4}{13.6}{-12.1}{alt}{[workload \GT\ threshold]}{8.5}
\draw[dashed,thick,black] (6.5,-9.8) -- (13.6,-9.8);
\node[anchor=north west,font=\footnotesize] at (6.6,-9.8) {[else]};
\frag{6.5}{-17.5}{17.6}{-21.3}{opt}{[idle \GT\ timeout]}{8.5}

\act{\xI}{-1.5}{-2.7} \act{\xI}{-15.6}{-16.7}
\act{\xK}{-2.7}{-13.2} \act{\xK}{-18.6}{-20.2}
\act{\xO}{-3.9}{-6.5} \act{\xO}{-14.2}{-14.7}
\act{\xP}{-13.2}{-20.7}

\sq{sm}{\xD}{\xI}{-1.5}{10: click ``Run''}
\sq{sm}{\xI}{\xK}{-2.7}{11: runRequest(roomId, ImageSpec)}
\sq{sm}{\xK}{\xO}{-3.9}{12: evaluateDemand(metrics)}
\selfm{\xO}{-5.0}{13: runPSO() / runACO()}{3.0}
\sq{sr}{\xO}{\xK}{-6.5}{14: ResourcePlan (node, cpu, freq)}

\selfm{\xK}{-8.5}{15a: scaleOut(): replicate pods (HPA / ARAA)}{2.9}
\selfm{\xK}{-10.9}{15b: reuse warm pod}{2.9}

\sql{sm}{\xK}{\xP}{-13.2}{16: deployContainer(spec, plan)}{0.5}
\sq{sm}{\xO}{\xP}{-14.4}{17: applyDVFS(freq)}
\sql{sr}{\xP}{\xI}{-15.6}{18: streamLogs(stdout, stderr)}{0.5}
\sq{sr}{\xI}{\xD}{-16.7}{19: render console output}

\selfm{\xK}{-18.6}{20: monitorIdle() (ARKLM)}{2.9}
\sq{sm}{\xK}{\xP}{-20.2}{21: terminateIdlePod(podId)}
\draw[very thick,black] (\xP-0.3,-20.5) -- (\xP+0.3,-21.1) (\xP-0.3,-21.1) -- (\xP+0.3,-20.5);
\end{tikzpicture}}
\caption{Sequence Diagram, Part B: Execution, Scaling, Log Streaming, and Teardown}
\label{fig:seq_b}
\end{figure}

Figure~\ref{fig:seq_b} illustrates the workflow for code execution, scaling, and session teardown. When a user runs code, the system validates current utilization, executes PSO/ACO optimization, updates container limits in place, modulates DVFS frequency, streams console logs to all connected clients, and triggers idle snapshotting when inactivity timeouts expire.

\clearpage
\subsection{State Machine Diagram: Workspace Pod Lifecycle}
\begin{figure}[H]
\centering
\singlespacing
\pngdiag{state_diag.png}{%
\begin{tikzpicture}[font=\footnotesize, xscale=1.0, yscale=1.1,
  st/.style={rectangle split, rectangle split parts=2, thick, rounded corners=2pt,
             text=black, text width=3.5cm, align=center, inner sep=3pt},
  tr/.style={-{Straight Barb[length=2.6mm,width=2.8mm]}, thick, black}]

\draw[thick, rounded corners=8pt, black] (0,-7.2) rectangle (16.6,-17.4);
\node[anchor=north west,font=\small\bfseries] at (0.3,-7.3) {Active\ {\tiny(composite state)}};

\node[circle,draw=black,very thick,minimum size=0.6cm,inner sep=0] (i0) at (2,2.6) {};
\fill[black] (i0) circle (0.16);
\state{pend}{2,0}{Pending}{entry / await ImageSpec\\do / queue request}
\state{sch}{8,0}{Scheduling}{entry / query Optimizer\\do / run PSO or ACO}
\state{prov}{14,0}{Provisioning}{entry / create pod\\do / pull image}
\state{fail}{8,-4.4}{Failed}{entry / log error, alert\\exit / release partial resources}
\state{init}{14,-4.4}{Initializing}{do / clone repo\\do / install deps, start IDE}

\node[circle,draw=black,very thick,minimum size=0.55cm,inner sep=0] (i1) at (0.9,-10.8) {};
\fill[black] (i1) circle (0.14);
\state{col}{4.8,-10.8}{Collaborating}{do / sync deltas, presence\\entry / restore full frequency}
\state{exe}{12.2,-10.8}{Executing}{do / run build or tests\\do / stream logs}
\state{scl}{12.2,-14.6}{Scaling Out}{do / replicate pods\\(HPA / ARAA)}
\state{idle}{4.8,-15.4}{Idle / Power-Saving}{entry / DVFS to min frequency\\do / ARKLM idle timer}

\state{term}{4.8,-20.6}{Terminating}{do / save workspace snapshot\\exit / delete pod}
\node[circle,fill=black,minimum size=0.6cm,inner sep=0] (fin) at (4.8,-23.6) {};

\draw[tr] (i0) -- node[tl,right,xshift=1pt] {open room} (pend.north);
\draw[tr] (pend.east) -- node[tl,above,text width=2.3cm] {ImageSpec ready / queryOptimizer()} (sch.west);
\draw[tr] (sch.east) -- node[tl,above,text width=2.3cm] {ResourcePlan ready / pullImage()} (prov.west);
\draw[tr] ([xshift=-0.8cm]sch.north) .. controls +(-0.3,1.1) and +(0.3,1.1) .. ([xshift=0.8cm]sch.north);
\node[tl,above,text width=3.2cm] at ($(sch.north)+(0,0.85)$) {[no capacity] / backoff and retry};
\draw[tr] (prov.south) -- node[tl,right,text width=2.3cm] {container created / mountWorkspace()} (init.north);
\draw[tr] (init.west) -- node[tl,above,text width=2.2cm] {build error / timeout} (fail.east);
\draw[tr] (fail.north) -- node[tl,left,text width=1.8cm] {[retries \LT\ max]} (sch.south);
\draw[tr] (init.south) -- node[tl,right,text width=2.3cm] {pod ready / readinessProbe OK} (14,-7.2);
\draw[tr] (8,-7.2) -- node[tl,right,text width=2.2cm] {pod crash (OOM / node fail)} (fail.south);
\draw[tr] (fail.west) -- node[tl,above,pos=0.3,text width=2.4cm] {retries exhausted / cleanup} (-1.0,-4.4) -- (-1.0,-20.6) -- (term.west);

\draw[tr] (i1) -- (col.west);
\draw[tr] ([xshift=-0.8cm]col.north) .. controls +(-0.3,1.1) and +(0.3,1.1) .. ([xshift=0.8cm]col.north);
\node[tl,above,text width=3.2cm] at ($(col.north)+(0,0.85)$) {delta received / applyOperation()};
\draw[tr] ([yshift=3.5mm]col.east) -- node[tl,above,text width=3.3cm] {Run signal / execute()} ([yshift=3.5mm]exe.west);
\draw[tr] ([yshift=-3.5mm]exe.west) -- node[tl,below,text width=3.3cm] {Execution done / streamLogs()} ([yshift=-3.5mm]col.east);
\draw[tr] ([xshift=-9mm]exe.south) -- node[tl,left,text width=2.3cm] {[load \GT\ limit] / scaleOut()} ([xshift=-9mm]scl.north);
\draw[tr] ([xshift=9mm]scl.north) -- node[tl,right,text width=2.0cm] {replicas ready} ([xshift=9mm]exe.south);
\draw[tr] ([xshift=-9mm]col.south) -- node[tl,left,text width=2.4cm] {no input \GT\ T\_idle / lowerFrequency()} ([xshift=-9mm]idle.north);
\draw[tr] ([xshift=9mm]idle.north) -- node[tl,right,text width=2.2cm] {user input / restoreFrequency()} ([xshift=9mm]col.south);

\draw[tr] (idle.south) -- node[tl,right,text width=2.4cm] {session timeout (ARKLM)} (term.north);
\draw[tr] (10,-17.4) |- node[tl,above,pos=0.8,text width=2.8cm] {all users left / admin evict} (term.east);
\draw[tr] (term.south) -- node[tl,right,text width=1.8cm] {resources released} (fin.north);
\end{tikzpicture}}
\caption{State Machine Diagram: Workspace Pod Lifecycle}
\label{fig:statemachine}
\end{figure}

Figure~\ref{fig:statemachine} tracks the lifecycle states of a workspace container. Starting from `Pending', the pod transitions through `Scheduling', `Provisioning', and `Initializing' before entering the composite `Active' state. Within `Active', the pod transitions between `Collaborating', `Executing', `Scaling Out', and `Idle / Power-Saving'. When inactivity limits are reached, the pod enters `Terminating', persisting workspace state before releasing compute resources.

\clearpage
\subsection{Activity Diagram: Optimization and Allocation}
\begin{figure}[H]
\centering
\singlespacing
\pngdiag{activity_diag.png}{%
\begin{tikzpicture}[font=\footnotesize, xscale=1.05, yscale=1.05,
  act/.style={rectangle, draw=black, thick, fill=white, text=black, text width=3.1cm, align=center,
              minimum height=0.85cm, inner sep=3pt},
  dec/.style={diamond, aspect=2.2, draw=black, thick, fill=white, text=black, text width=1.6cm, align=center,
              inner sep=0pt, font=\scriptsize},
  fl/.style={-{Straight Barb[length=2.6mm,width=2.8mm]},thick,black},
  fb/.style={font=\scriptsize, fill=white, inner sep=1pt},
  bar/.style={rectangle, fill=black, minimum height=0.14cm, inner sep=0},
  note/.style={rectangle, fill=white, draw=black, text=black, font=\scriptsize, align=center,
               inner sep=3pt, minimum width=1.9cm}]

\foreach \cx/\tx in {2.2/{Developer / IDE},6.6/{K8s Orchestrator},11.0/{Optimizer (PSO / ACO / DVFS)},15.4/{Worker Pod}}{
  \draw[thick,black] (\cx-2.2,0) rectangle (\cx+2.2,-27.9);
  \node[draw=black,thick,fill=white,minimum width=4.4cm,minimum height=0.8cm,font=\small\bfseries,text width=4.1cm,align=center] at (\cx,-0.4) {\tx};
}
\node[circle,draw=black,fill=black,minimum size=0.4cm,inner sep=0] (st) at (2.2,-1.5) {};
\node[act] (d1) at (2.2,-2.6) {Click Run / Build};
\node[act] (d2) at (2.2,-4.0) {Send runRequest\\(roomId, ImageSpec)};

\node[act] (o1) at (6.6,-5.4) {Receive and validate request};
\node[act] (o2) at (6.6,-6.8) {Collect metrics: CPU, RAM, queue length};
\node[dec] (dc) at (6.6,-8.6) {Workload \GT\ limit?};

\node[act] (aco) at (11.0,-8.6)  {Run ACO: schedule tasks on existing nodes};
\node[act] (dvl) at (11.0,-10.3) {Lower DVFS frequency (save energy)};
\node[act] (pso) at (11.0,-12.0) {Run PSO: size CPU / RAM / replicas};

\node[act] (o3) at (6.6,-13.6) {Trigger pod replication (HPA / ARAA)};
\node[diamond,draw=black,thick,fill=white,minimum width=0.8cm,minimum height=0.8cm,inner sep=0] (mg) at (6.6,-15.0) {};
\node[act] (o4) at (6.6,-16.4) {Provision / select worker pod};

\node[act] (p1) at (15.4,-16.4) {Pull image and mount workspace};
\node[act] (p2) at (15.4,-17.9) {Execute code in container (resource limits)};
\node[bar,minimum width=3.6cm] (fk) at (15.4,-19.1) {};
\node[act,text width=1.45cm] (p3) at (14.3,-20.5) {Stream logs to room};
\node[act,text width=1.45cm] (p4) at (16.5,-20.5) {Report CPU / RAM / energy};
\node[bar,minimum width=3.6cm] (jn) at (15.4,-21.9) {};

\node[act] (d3) at (2.2,-23.2) {Render output in all consoles};
\node[act] (o5) at (6.6,-23.2) {Start idle timer (ARKLM)};
\node[dec] (dc2) at (6.6,-25.0) {Idle \GT\ timeout?};
\node[act] (keep) at (11.0,-25.0) {Keep pod alive; DVFS power-saving mode};
\node[act] (o6) at (6.6,-26.7) {Terminate pod and release resources};
\node[circle,draw=black,very thick,minimum size=0.6cm,inner sep=0] (fin) at (6.6,-27.55) {};
\fill[black] (fin) circle (0.17);

\node[note] (n1) at (-2.4,-1.5) {Developer is\\in a room};
\draw[dashed,thick,black] (n1.east) -- (st.west);
\node[note] (n2) at (11.0,-27.2) {Pod released,\\resources free};
\draw[dashed,thick,black] (fin.east) -- (n2.west);

\draw[fl] (st) -- (d1);
\draw[fl] (d1) -- (d2);
\draw[fl] (d2.south) |- (o1.west);
\draw[fl] (o1) -- (o2);
\draw[fl] (o2) -- (dc);
\draw[fl] (dc.east) -- node[fb,above] {[No]} (aco.west);
\draw[fl] (aco) -- (dvl);
\draw[fl] (dc.south) |- node[fb,pos=0.12,right] {[Yes]} (pso.west);
\draw[fl] (dvl.east) -- ++(0.35,0) |- (mg.east);
\draw[fl] (pso.south) |- (o3.east);
\draw[fl] (o3) -- (mg);
\draw[fl] (mg) -- (o4);
\draw[fl] (o4) -- (p1);
\draw[fl] (p1) -- (p2);
\draw[fl] (p2) -- (fk);
\draw[fl] ([xshift=-1.1cm]fk.south) -- (p3.north);
\draw[fl] ([xshift=1.1cm]fk.south) -- (p4.north);
\draw[fl] (p3.south) -- ([xshift=-1.1cm]jn.north);
\draw[fl] (p4.south) -- ([xshift=1.1cm]jn.north);
\draw[fl] (jn.west) -| (d3.north);
\draw[fl] (d3) -- (o5);
\draw[fl] (o5) -- (dc2);
\draw[fl] (dc2.east) -- node[fb,above] {[No]} (keep.west);
\draw[fl] (keep.north) |- node[fb,pos=0.78,above] {re-check} (o5.east);
\draw[fl] (dc2.south) -- node[fb,right] {[Yes]} (o6);
\draw[fl] (o6) -- (fin);
\end{tikzpicture}}
\caption{Activity Diagram: Resource Allocation and Execution Flow}
\label{fig:activity}
\end{figure}

Figure~\ref{fig:activity} presents an activity diagram organized into four swimlanes: Developer/IDE, Kubernetes Orchestrator, Optimizer (PSO/ACO/DVFS), and Worker Pod. It maps the decision logic for workload evaluation, directing normal interactive loads to ACO node balancing and DVFS power saving, and routing heavy compute bursts to PSO vertical scaling and pod replication.

\clearpage
\subsection{Component Diagram}
\begin{figure}[H]
\centering
\singlespacing
\pngonly{comp_diag.png}
\caption{Component Diagram (Provided and Required Interfaces)}
\label{fig:component}
\end{figure}

Figure~\ref{fig:component} illustrates the major software components and their interface dependencies. The frontend components (\texttt{Monaco Editor Component}, \texttt{Terminal Component}, and \texttt{CRDT Client}) communicate with edge ingress services through WebSocket and REST interfaces, which in turn connect to the collaboration engine, execution controllers, and telemetry scrapers.

\clearpage
\subsection{Deployment Diagram}
\begin{figure}[H]
\centering
\singlespacing
\pngonly{deploy_diag.png}
\caption{Deployment Diagram: Physical and Virtual Topologies}
\label{fig:deployment}
\end{figure}

Figure~\ref{fig:deployment} models the physical deployment topology. It displays the master control-plane node and bare-metal worker nodes running under containerd and cgroups v2, detailing network ingress routing, container network interface (CNI) bridging, and Prometheus metric scraping endpoints.

\clearpage
\subsection{Object Diagram: Runtime Instance Snapshot}
\begin{figure}[H]
\centering
\singlespacing
\pngdiag{object_diag.png}{%
\begin{tikzpicture}[font=\small, xscale=1.0, yscale=1.1,
  obj/.style={rectangle split, rectangle split parts=2, draw=black, thick,
              rectangle split part fill={white,white},
              text width=6.2cm, align=left, inner sep=4pt}]

\inst{alice}{0,0}{devAlice : UserSession}{userId = ``usr\_9801''\\cursorOffset = 142\\role = COLLABORATOR}
\inst{bob}{10,0}{devBob : UserSession}{userId = ``usr\_9802''\\cursorOffset = 188\\role = COLLABORATOR}
\inst{room}{5,-5.0}{activeRoom : CollabRoom}{roomId = ``room\_c109''\\activeParticipants = 2\\currentRevision = 894}
\inst{doc}{0,-10.2}{mainBuffer : DocumentBuffer}{fileUri = ``/src/index.ts''\\lengthBytes = 14209\\pendingOps = 0}
\inst{plan}{10,-10.2}{currentPlan : ResourcePlan}{targetNode = ``node-worker-03''\\allocatedCpu = 1.50 cores\\allocatedRamMB = 2048\\targetFreqMHz = 2400}
\inst{pod}{5,-15.6}{runnerPod : WorkspacePod}{podId = ``pod-c109-exec-7''\\containerStatus = ``RUNNING''\\lastHeartbeat = 1728389100\\governorMode = ``performance''}

\draw[thick,black] (alice.south) -- (room.north west) node[lbl,pos=0.5,left=1pt] {connectedTo};
\draw[thick,black] (bob.south) -- (room.north east) node[lbl,pos=0.5,right=1pt] {connectedTo};
\draw[thick,black] (room.south west) -- (doc.north) node[lbl,pos=0.5,left=1pt] {locks and manages};
\draw[thick,black] (room.south east) -- (plan.north) node[lbl,pos=0.5,right=1pt] {evaluates};
\draw[thick,black] (plan.south) -- (pod.north east) node[lbl,pos=0.5,right=1pt] {configures};
\draw[thick,black] (doc.south) -- (pod.north west) node[lbl,pos=0.5,left=1pt] {mounted inside};
\end{tikzpicture}}
\caption{Object Diagram: Runtime Instance Snapshot}
\label{fig:object}
\end{figure}

Figure~\ref{fig:object} provides a concrete snapshot of system objects during active execution. It illustrates two collaborating users (\texttt{devAlice} and \texttt{devBob}) sharing room \texttt{room\_c109} and editing \texttt{/src/index.ts}. The optimizer has calculated a \texttt{ResourcePlan} of 1.50 CPU cores, 2048 MB RAM, and 2400 MHz clock frequency for container \texttt{pod-c109-exec-7}.

\clearpage
\subsection{Collaboration (Communication) Diagram}
\begin{figure}[H]
\centering
\singlespacing
\pngonly{colab_diag.png}
\caption{Collaboration (Communication) Diagram}
\label{fig:communication}
\end{figure}

Figure~\ref{fig:communication} illustrates the sequenced message paths between runtime components during active editing and execution, documenting the step-by-step coordination between client browsers, WebSocket brokers, the CRDT engine, and the Kubernetes pod controller.

\clearpage
\subsection{Package Diagram}
\begin{figure}[H]
\centering
\singlespacing
\pngonly{package_diag.png}
\caption{Package Diagram: Module Dependencies}
\label{fig:package}
\end{figure}

Figure~\ref{fig:package} organizes the codebase into functional packages: \texttt{ClientCore}, \texttt{CollabEngine}, \texttt{InfrastructureManager}, \texttt{OptimizationSuite}, and \texttt{CommonTelemetry}. Layered dependency rules prevent cyclic coupling across system boundaries.

\clearpage
\subsection{Composite Structure Diagram}
\begin{figure}[H]
\centering
\singlespacing
\pngonly{Composite_diag.png}
\caption{Composite Structure Diagram: Internal Workspace Session}
\label{fig:composite}
\end{figure}

Figure~\ref{fig:composite} shows the internal structure of an active \texttt{WorkspaceSession} component, detailing its internal communication ports, CRDT synchronization caches, terminal multiplexers, and container execution adapters.

% =========================================================
% DATA DESIGN
% =========================================================
\clearpage
\section{Data Design and Database Schema}

\subsection{Internal Software Data Structures}
In the collaboration runtime, the CRDT document sequence is maintained in memory using a combination of contiguous arrays and balance-indexed skip lists. Each character atom contains:
\begin{itemize}
    \item \texttt{id}: Composite identifier \texttt{(client\_id, clock\_tick)}.
    \item \texttt{origin\_left} and \texttt{origin\_right}: Boundary anchor identifiers for deterministic insertion.
    \item \texttt{content}: String atom payload.
    \item \texttt{deleted}: Boolean tombstone flag.
\end{itemize}
Collaborator states are tracked in an in-memory hash map (\texttt{userId} $\rightarrow$ \texttt{\{cursorLine, cursorCol, colorHex, lastHeartbeat\}}).

\subsection{Global Data Structures}
Horizontal scaling across multiple collaboration worker pods uses Redis Pub/Sub channels. Channel naming conventions follow \texttt{room:<roomId>:events} for delta synchronization and presence heartbeats. A shared Redis sorted set tracks room activity timestamps to drive the ARKLM eviction engine.

\subsection{Temporary Data Structures}
Prometheus and cAdvisor maintain 60-second in-memory ring buffers containing sliding-window telemetry for each active pod:
\begin{equation}
    \mathcal{M}_{\text{pod}} = \{\langle t_k, \text{cpu\_usage}_k, \text{mem\_rss}_k, \text{io\_ops}_k \rangle \mid k \in [1, 60]\}
\end{equation}

\subsection{Database Schema Design}
Persistent entities are structured in PostgreSQL:
\begin{itemize}
    \item \textbf{users:} \texttt{(id, username, email, password\_hash, oauth\_provider, created\_at)}
    \item \textbf{workspaces:} \texttt{(id, owner\_id, repo\_url, branch, last\_snapshot\_uri, created\_at)}
    \item \textbf{collab\_rooms:} \texttt{(id, workspace\_id, active\_participants, room\_status, updated\_at)}
    \item \textbf{sla\_audit\_logs:} \texttt{(id, pod\_id, timestamp, cpu\_demand, cpu\_allocated, memory\_demand, memory\_allocated, violation\_flag)}
\end{itemize}

% =========================================================
% PROTOTYPING & BENCHMARKS
% =========================================================
\section{Prototyping and Experimental Validation}
A fully functional prototype was deployed and evaluated across a multi-node Kubernetes cluster to benchmark collaborative synchronization latency, container scaling, and power consumption.

\begin{figure}[H]
    \centering
    \safeincludegraphics[width=0.85\linewidth]{7.jpeg}{CollabIDE Authentication Interface}
    \caption{CollabIDE Authentication Interface: Clean Login/OAuth and Live Session Preview}
    \label{fig:prototype_auth}
\end{figure}

Figure~\ref{fig:prototype_auth} illustrates the authentication interface where developers log in via OAuth2 credentials and view active workspace rooms and current participant counts before joining.

\begin{figure}[H]
    \centering
    \safeincludegraphics[width=0.88\linewidth]{1.jpeg}{Interactive Microservice Topologies}
    \caption{Interactive Microservice Topologies: Node.js and Python Runtime Sizing Canvas}
    \label{fig:prototype_canvas}
\end{figure}

Figure~\ref{fig:prototype_canvas} shows the workspace configuration canvas. Developers can parameterize runtime environments (Python 3.11, Node.js 20) with baseline CPU quotas (0.50 cores), memory limits (256 MB), and autoscaling bounds before pod creation.

\begin{figure}[H]
    \centering
    \safeincludegraphics[width=0.90\linewidth]{2.jpeg}{Real-Time Resource Optimizer Dashboard}
    \caption{Real-Time Resource Optimizer Dashboard: Live Power Curve and Policy Evaluation}
    \label{fig:prototype_optimizer}
\end{figure}

Figure~\ref{fig:prototype_optimizer} displays the live closed-loop optimizer dashboard. It evaluates candidate resource allocation policies in real time, plotting host power draw, utilization curves, and SLA violation metrics under variable developer workloads.

\begin{table}[htbp]
\centering
\caption{Real-Time Optimization Strategy Evaluation (Measured Prototype Output)}
\label{tab:prototype_metrics}
\begin{tabularx}{0.85\linewidth}{@{}lXXX@{}}
\toprule
\textbf{Allocation Strategy} & \textbf{Host Power} & \textbf{CPU Util} & \textbf{Unmet SLA} \\ \midrule
\textbf{DVFS Only (Proposed)} & \textbf{27.4 W (lowest)} & \textbf{60\%} & \textbf{0.0\%} \\
Static Allocation & 105.4 W & 6\% & 0.0\% \\
K8s-style Least-Requested & 105.4 W & 6\% & 0.0\% \\
ACO Placement & 73.4 W & 12\% & 0.0\% \\ \bottomrule
\end{tabularx}
\end{table}

As shown in Table~\ref{tab:prototype_metrics}, combining DVFS with adaptive container rightsizing reduced host power draw from 105.4~W down to 27.4~W---a 74\% energy saving---while maintaining zero SLA violations throughout the benchmark period.

% =========================================================
% CHAPTER 6: PROJECT PLANNING
% =========================================================
\chapter{PROJECT PLANNING}

\section{Project Timeline, Milestones, and Gantt Chart}
The engineering lifecycle for the Resource-Aware Collaborative Cloud IDE is structured across a multi-semester schedule spanning academic year 2026--2027. The project is divided into two primary execution stages: Stage-I (Semester I, July 2026 to November 2026) focusing on theoretical formulations, architecture design, core microservice scaffolding, and prototype benchmarking; and Stage-II (Semester II, December 2026 to April 2027) addressing production deployment, deep learning workload forecasting, kernel-level live container migration, and comprehensive multi-node stress testing.

Table~\ref{tab:detailed_schedule} details the monthly milestones, deliverables, and team ownership assignments across the 10-month development cycle.

\begin{table}[htbp]
\centering
\scriptsize
\caption{Detailed Monthly Project Schedule and Milestone Breakdown}
\label{tab:detailed_schedule}
\begin{tabularx}{\linewidth}{@{}llllX@{}}
\toprule
\textbf{Month} & \textbf{Phase} & \textbf{Lead} & \textbf{Status} & \textbf{Core Deliverables and Engineering Milestones} \\ \midrule
\textbf{Jul 2026} & Literature \& Feasibility & Bikash & Completed & Comprehensive survey on sequence CRDTs vs OT; literature review on swarm metaheuristics (PSO, ACO) and kernel DVFS. \\
\textbf{Aug 2026} & SRS \& Architecture & Abhay & Completed & Complete Software Requirements Specification (SRS); design of multi-layer microservice architecture; API specifications. \\
\textbf{Sep 2026} & Bare-Metal Cluster Setup & Anish P. & Completed & Assembly of 3-node bare-metal Kubernetes v1.28 cluster; containerd configuration; cAdvisor sub-second scraping setup. \\
\textbf{Oct 2026} & Collaborative Engine & Bikash & Completed & Yjs sequence CRDT integration; WebSocket gateway; Redis Pub/Sub room bus; sub-100ms peer synchronization validation. \\
\textbf{Nov 2026} & Optimization Prototype & Anish K. & Completed & Implementation of PSO vertical rightsizing; ACO pod scheduling; host DVFS daemon; Stage-I preliminary report delivery. \\ \midrule
\textbf{Dec 2026} & Stage-II Inception & Abhay & Planned & Refactoring optimizer core into native C++/Go extensions; integrating Prometheus Operator for enterprise metric reliability. \\
\textbf{Jan 2027} & Predictive Modeling & Anish P. & Planned & Implementation of LSTM/Transformer time-series models for proactive developer burst forecasting from keystroke metrics. \\
\textbf{Feb 2027} & In-Place Pod Migration & Bikash & Planned & Integration of Checkpoint/Restore in Userspace (CRIU) for zero-downtime, stateful live container migration during rebalancing. \\
\textbf{Mar 2027} & Multi-Tenant Stress Test & Anish K. & Planned & Large-scale multi-tenant load testing (100+ concurrent rooms); synthetic compilation burst injection; energy audits. \\
\textbf{Apr 2027} & Final Documentation & All & Planned & Final Stage-II project report submission, academic paper publication, open-source code release, and defense presentation. \\ \bottomrule
\end{tabularx}
\end{table}

\subsection{Visual Project Schedule (Gantt Chart)}
Figure~\ref{fig:gantt_chart} illustrates the monthly distribution of project tasks, depicting critical paths, concurrent engineering phases, and milestone dependencies across the 10-month project timeline.

\begin{figure}[H]
\centering
\singlespacing
\pngdiag{gantt_chart.png}{%
\begin{tikzpicture}[xscale=1.1, yscale=0.75, font=\scriptsize]
    \draw[help lines, color=gray!30] (0,0) grid (10,10);
    
    \foreach \m/\x in {Jul/0.5, Aug/1.5, Sep/2.5, Oct/3.5, Nov/4.5, Dec/5.5, Jan/6.5, Feb/7.5, Mar/8.5, Apr/9.5} {
        \node[font=\bfseries] at (\x, 10.5) {\m};
    }
    
    \filldraw[rounded corners=2pt, fill=blue!40, draw=blue!80!black, thick] (0, 8.65) rectangle (1.5, 9.35);
    \node[anchor=west, font=\scriptsize\bfseries] at (0.1, 9) {1. Literature Survey \& Feasibility};

    \filldraw[rounded corners=2pt, fill=blue!40, draw=blue!80!black, thick] (1, 7.65) rectangle (2.5, 8.35);
    \node[anchor=west, font=\scriptsize\bfseries] at (0.1, 8) {2. Requirements \& Architecture Design};

    \filldraw[rounded corners=2pt, fill=blue!40, draw=blue!80!black, thick] (2, 6.65) rectangle (3.5, 7.35);
    \node[anchor=west, font=\scriptsize\bfseries] at (0.1, 7) {3. K8s Infrastructure \& cAdvisor Telemetry};

    \filldraw[rounded corners=2pt, fill=blue!40, draw=blue!80!black, thick] (3, 5.65) rectangle (4.5, 6.35);
    \node[anchor=west, font=\scriptsize\bfseries] at (0.1, 6) {4. CRDT Collaboration \& WebSocket Engine};

    \filldraw[rounded corners=2pt, fill=blue!40, draw=blue!80!black, thick] (3.8, 4.65) rectangle (5.0, 5.35);
    \node[anchor=west, font=\scriptsize\bfseries] at (0.1, 5) {5. PSO/ACO/DVFS Optimization Prototype};
    
    \draw[thick, dashed, red] (5.0, 0) -- (5.0, 10);
    \node[anchor=south, color=red!80!black, font=\tiny\bfseries, rotate=90] at (5.1, 2.5) {Stage-I Review (Nov 2026)};
    
    \filldraw[rounded corners=2pt, fill=teal!40, draw=teal!80!black, thick] (5.0, 3.65) rectangle (6.2, 4.35);
    \node[anchor=west, font=\scriptsize\bfseries] at (0.1, 4) {6. Optimizer Refactoring \& Prom-Operator};

    \filldraw[rounded corners=2pt, fill=teal!40, draw=teal!80!black, thick] (6.0, 2.65) rectangle (7.5, 3.35);
    \node[anchor=west, font=\scriptsize\bfseries] at (0.1, 3) {7. Predictive LSTM Workload Modeling};

    \filldraw[rounded corners=2pt, fill=teal!40, draw=teal!80!black, thick] (7.0, 1.65) rectangle (8.5, 2.35);
    \node[anchor=west, font=\scriptsize\bfseries] at (0.1, 2) {8. CRIU Stateful Pod Live Migration};

    \filldraw[rounded corners=2pt, fill=teal!40, draw=teal!80!black, thick] (8.0, 0.65) rectangle (9.5, 1.35);
    \node[anchor=west, font=\scriptsize\bfseries] at (0.1, 1) {9. Multi-Tenant Stress Testing \& Audits};

    \filldraw[rounded corners=2pt, fill=teal!40, draw=teal!80!black, thick] (8.8, -0.35) rectangle (10.0, 0.35);
    \node[anchor=west, font=\scriptsize\bfseries] at (0.1, 0) {10. Stage-II Documentation \& Defense};
\end{tikzpicture}}
\caption{Detailed 10-Month Project Gantt Chart (Stage-I and Stage-II Schedule)}
\label{fig:gantt_chart}
\end{figure}

\section{Technology Stack, Estimates, and Resource Reuse}
Table~\ref{tab:stack} summarizes the technology stack, software engineering estimates, and component reuse strategy for the platform.

\begin{table}[H]
\centering
\caption{Technology Stack and Project Estimates}
\label{tab:stack}
\begin{tabularx}{\linewidth}{@{}p{3.5cm}X@{}}
\toprule
\textbf{Item} & \textbf{Technical Specification} \\ \midrule
Front End & React 18, Monaco Editor engine, xterm.js terminal, Yjs CRDT client \\
Back End & Node.js 20 (y-websocket, REST services), Python 3.11 (PSO/ACO optimizers) \\
Persistence & PostgreSQL 16 (relational data), Redis 7.2 (Pub/Sub message bus) \\
Infrastructure & Docker, containerd, Kubernetes v1.28, Prometheus, cAdvisor \\
Hardware Testbed & 1 Control-plane node, 2 Worker nodes (Intel Core i7/Xeon, ACPI cpufreq enabled) \\
Code Size Estimate & $\approx 8,500$--$11,000$ lines of TypeScript and Python across $\sim 42$ classes \\
Component Reuse & Yjs, Monaco, xterm.js, and Prometheus are reused via open-source APIs; optimization algorithms are developed custom in-house \\
Project Completion & Stage-I completed; Stage-II production deployment scheduled for Semester II, 2026-27 \\ \bottomrule
\end{tabularx}
\end{table}

% =========================================================
% CHAPTER 7: SYSTEM TESTING AND VALIDATION
% =========================================================
\chapter{SYSTEM TESTING AND VALIDATION}

\section{Testing Strategy and Methodology}
Validating a distributed, resource-aware cloud development environment demands a multi-tiered verification methodology spanning unit verification, asynchronous integration testing, SLA stress injection, and long-term stability audits under high concurrency. The testing strategy isolates the verification process into four functional vectors:
\begin{enumerate}
    \item \textbf{CRDT Eventual Consistency and Convergence Testing:} Simulating synthetic network partitions, dropped WebSocket packets, out-of-order delta receipts, and high-frequency concurrent typing bursts to prove zero document divergence across connected peers.
    \item \textbf{Sub-Second Telemetry Ingestion Precision:} Verifying that cAdvisor and Prometheus pipelines accurately reflect physical container utilization without metric skew or scrape dropouts under rapid compute bursts.
    \item \textbf{In-Place Scaling Responsiveness:} Benchmarking the end-to-end response time between detecting a compute burst and updating Linux cgroups v2 boundaries via the Kubernetes pod patch API.
    \item \textbf{Hardware DVFS Stability and Transition Overhead:} Measuring latency and voltage stabilization times when modulating physical core frequencies between $f_{\text{low}}$ and $f_{\max}$.
\end{enumerate}

\section{Test Cases and Results}
Table~\ref{tab:test_cases} details representative system test cases, outlining target verification vectors, execution conditions, expected behaviors, and validated outcomes.

\begin{table}[H]
\centering
\scriptsize
\caption{System Integration and Validation Test Suite}
\label{tab:test_cases}
\begin{tabularx}{\linewidth}{@{}p{1.6cm}p{2.8cm}p{3.4cm}Xp{1.2cm}@{}}
\toprule
\textbf{Test ID} & \textbf{Test Scenario} & \textbf{Input Conditions} & \textbf{Expected Result} & \textbf{Status} \\ \midrule
\textbf{TC-CRDT-01} & Concurrent Inserts & 10 clients insert characters at offset 0 simultaneously. & Deterministic ordering via Lamport ticks; zero dropped characters. & PASS \\
\textbf{TC-CRDT-02} & Network Partition & 2 clients disconnected for 60s; 100 offline edits made. & Clean bidirectional merge upon reconnection; buffers converge. & PASS \\
\textbf{TC-OPT-01} & Idle Power Down & No editor keystrokes or terminal I/O for 30 seconds. & Host CPU governor drops frequency to $f_{\text{low}}$; power drops. & PASS \\
\textbf{TC-OPT-02} & Compilation Burst & User runs GCC multi-core build; CPU demand spikes to 100\%. & PSO expands quota to 2.0 cores in $<1.2$s; no OOM or throttle. & PASS \\
\textbf{TC-LIF-01} & ARKLM Eviction & Room remains inactive for 15 minutes ($T_{\text{idle}} > 900$s). & Container snapshotted to S3/Ceph; pod deleted; memory freed. & PASS \\
\textbf{TC-SEC-01} & Sandbox Escape & Process attempts root escalation and raw socket bind. & Blocked by seccomp filter and cgroups v2; alert triggered. & PASS \\ \bottomrule
\end{tabularx}
\end{table}

\Needspace{10\baselineskip}
\section{Performance Benchmarking and Load Testing}
To evaluate scalability, synthetic workloads were injected into the cluster using automated headless browser clients. Peer editing propagation latency was tracked across increasing room concurrency levels.

\begin{table}[H]
\centering
\small
\caption{Peer Edit Synchronization Latency vs. Concurrent Collaborators}
\label{tab:latency_benchmarks}
\begin{tabularx}{\linewidth}{@{}XXXXX@{}}
\toprule
\textbf{Concurrent Users} & \textbf{Mean Latency (ms)} & \textbf{95th Pct (ms)} & \textbf{99th Pct (ms)} & \textbf{CPU Util (\%)} \\ \midrule
2 Users & 14.2 & 18.5 & 22.1 & 1.2\% \\
5 Users & 21.6 & 28.4 & 34.2 & 2.8\% \\
10 Users & 42.1 & 58.7 & 76.4 & 6.4\% \\
15 Users & 68.4 & 89.2 & 112.5 & 11.8\% \\
20 Users & 94.7 & 124.6 & 158.2 & 18.5\% \\ \bottomrule
\end{tabularx}
\end{table}

As documented in Table~\ref{tab:latency_benchmarks}, the platform maintains sub-100ms peer synchronization latency up to 15 concurrent developers editing the same document buffer, validating the low overhead of the decentralized CRDT engine.

% =========================================================
% CHAPTER 8: SECURITY & THREAT MODELING
% =========================================================
\chapter{SECURITY, ISOLATION, AND THREAT MODELING}

\section{Container Sandbox Isolation Mechanisms}
Multi-tenant Cloud IDEs operate under an adversarial threat model: users are granted arbitrary shell execution access within their designated workspaces. To prevent cross-tenant data leakage, privilege escalation, and host-level denial of service, the execution engine enforces strict defense-in-depth isolation across three layers:
\begin{enumerate}
    \item \textbf{Linux Control Groups (cgroups v2):} Every workspace pod is placed inside a dedicated cgroup hierarchy. Limits are strictly enforced for maximum CPU core consumption (\texttt{cpu.max}), memory footprint (\texttt{memory.max}), swap allocation (\texttt{memory.swap.max}), and maximum task threads (\texttt{pids.max}). Restricting thread creation prevents fork-bomb attacks from starving host resources.
    \item \textbf{Linux Namespaces:} Workspaces run within isolated Mount, PID, Network, IPC, UTS, and User namespaces. Processes inside a pod cannot observe host processes or network sockets belonging to neighboring tenants.
    \item \textbf{Seccomp and AppArmor Profiles:} Default Docker seccomp profiles are tightened to block dangerous system calls, including \texttt{ptrace}, \texttt{sys\_admin}, \texttt{mount}, \texttt{reboot}, and raw socket creation primitives.
\end{enumerate}

\section{Network Policy and Ingress Hardening}
Container network interfaces (CNI) enforce isolated Calico network policies. Workspaces can access external Git repositories over HTTPS and download package dependencies from approved registries (e.g., npm, PyPI). However, all direct ingress connections to the container are blocked. Client terminal traffic is routed through authenticated WebSocket streams mediated by the platform API Gateway.

% =========================================================
% CHAPTER 9: CONCLUSIONS AND FUTURE WORK
% =========================================================
\chapter{CONCLUSIONS AND FUTURE WORK}

\section{Summary of Research Contributions}
This preliminary project report presented the comprehensive architectural design, theoretical modeling, algorithmic formulation, and prototype evaluation of a Resource-Aware Collaborative Cloud Integrated Development Environment. By establishing an intelligent, bidirectional closed-loop relationship between collaborative frontend editing actions and backend infrastructure scaling, the platform effectively resolves the fundamental tension between high-concurrency developer collaboration and dynamic cloud infrastructure efficiency.

At the collaboration layer, replacing centralized Operational Transformation (OT) with decentralized Conflict-free Replicated Data Types (CRDTs) provides mathematically verified eventual consistency with sub-100ms peer synchronization latency and zero central locking overhead. Concurrently, at the infrastructure layer, the integration of sub-second cAdvisor telemetry with Particle Swarm Optimization (PSO) enables continuous, in-place vertical rightsizing of container CPU quotas and memory boundaries without restarting pods. Complementing this, Ant Colony Optimization (ACO) balances pod allocations across worker nodes, Dynamic Voltage and Frequency Scaling (DVFS) modulates processor frequencies during developer idle periods, and the Adaptive Resource \& Knowledge Lifecycle Manager (ARKLM) snapshots dormant workspaces.

Empirical evaluations across multi-tenant cluster testbeds confirm that the proposed system lowers physical host power dissipation by up to 74\%, eliminates static over-provisioning, avoids SLA violations, and delivers responsive, conflict-free collaboration under variable multi-tenant developer workloads.

\section{Limitations}
\begin{enumerate}
    \item \textbf{Metric Ingestion Granularity:} Prometheus scrapes container metrics at 1-second intervals, introducing a brief observation lag before sharp compilation bursts trigger PSO scaling.
    \item \textbf{Hardware DVFS Constraints:} Direct CPU frequency modulation requires low-level ACPI cpufreq kernel access, which is readily supported on bare-metal or virtualization-passthrough nodes but restricted on fully managed public cloud container services.
    \item \textbf{Initial Container Warm-Up:} First-time workspace launches require pulling base images and mounting storage layers, causing cold-start latencies of 3 to 5 seconds without pre-warmed pod pools.
\end{enumerate}

\section{Future Scope}
Future enhancements planned for Stage-II include:
\begin{enumerate}
    \item \textbf{Deep Learning Workload Prediction:} Integrating LSTM or Transformer models to forecast developer compilation bursts from keystroke cadence and code syntax complexity.
    \item \textbf{WebAssembly Sandboxing:} Incorporating WebAssembly (WASM) micro-sandboxes alongside Docker containers for near-instantaneous execution of lightweight scripts.
    \item \textbf{Stateful Live Migration:} Implementing live container migration across heterogeneous physical cluster nodes using Checkpoint/Restore in Userspace (CRIU) to eliminate session interruptions during node rebalancing.
\end{enumerate}

% =========================================================
% CHAPTER 10: REFERENCES (COMPREHENSIVE BIBLIOGRAPHY)
% =========================================================
\chapter{REFERENCES}
{\renewcommand{\chapter}[2]{}%
\begin{thebibliography}{99}

\bibitem{borg}
A.~Verma, L.~Pedrosa, M.~Korupolu, D.~Oppenheimer, E.~Tune, and J.~Wilkes,
``Large-scale cluster management at Google with Borg,''
in \textit{Proc. 10th European Conf. on Computer Systems (EuroSys)}, pp. 1--17, 2015.

\bibitem{aras}
L.~Shan, Y.~Zhang, and X.~Wang,
``Adaptive resource allocation for workflow containerization on Kubernetes,''
\textit{Journal of Systems Architecture}, vol. 138, p. 102863, 2023.

\bibitem{pace}
A.~Al-Dhuraibi, M.~Paraiso, N.~Djarallah, and P.~Merle,
``Auto-scaling containerized cloud applications: A workload-driven approach,''
\textit{IEEE Transactions on Cloud Computing}, vol. 10, no. 4, pp. 2405--2420, 2022.

\bibitem{sata}
L.~Toka, G.~Dobreff, B.~Fodor, and B.~P.~Ger\H{o},
``SATA: The SLA-Adaptive Threshold Adjustment Algorithm for Kubernetes Horizontal Autoscaler,''
\textit{IEEE Transactions on Network and Service Management}, vol. 20, no. 2, pp. 1420--1433, 2023.

\bibitem{psoaco}
M.~Dubey and S.~Sharma,
``A hybrid multi-faceted task scheduling algorithm for cloud computing environment,''
\textit{Journal of Cloud Computing}, vol. 12, no. 1, pp. 45--58, 2023.

\bibitem{pso_lstm}
C.~H.~Hsu, S.~Chen, and M.~B.~Lin,
``KIF: A Kubernetes Pod Scheduling Model Based on Particle Swarm Optimization and Long Short-Term Memory Network,''
in \textit{Proc. IEEE Int. Conf. on Edge Computing (EDGE)}, pp. 88--97, 2024.

\bibitem{cshbo}
R.~S.~Rajesh, P.~Deepa, and K.~Suresh,
``Hybrid metaheuristic technique for optimal container resource allocation in cloud,''
\textit{Cluster Computing}, vol. 27, no. 3, pp. 3201--3218, 2024.

\bibitem{atsoa}
N.~A.~Al-Janabi and A.~K.~Mahmood,
``Adaptive Tree Seed Optimization Algorithm for Multi-Objective Task Scheduling in Cloud Computing,''
\textit{IEEE Access}, vol. 11, pp. 54201--54215, 2023.

\bibitem{tspra}
Z.~Wang, Y.~Ding, and H.~Sun,
``Temporal-aware semi-supervised performance prediction and resource allocation for cloud-native microservices,''
\textit{Future Generation Computer Systems}, vol. 158, pp. 312--326, 2024.

\bibitem{petas}
J.~Krasteva, R.~Dobreva, and G.~Georgiev,
``Efficient evolutionary optimization using predictive auto-scaling in containerized environment,''
\textit{Applied Soft Computing}, vol. 142, p. 110352, 2023.

\bibitem{cec}
Y.~Li, X.~Cheng, and Y.~Zhang,
``CEC: A Containerized Edge Computing Framework for Dynamic Resource Provisioning,''
\textit{IEEE Transactions on Services Computing}, vol. 16, no. 3, pp. 1890--1903, 2023.

\bibitem{dvfs}
K.~Beena, S.~Rathore, and M.~Patel,
``Adaptive energy optimization in cloud computing through containerization and DVFS,''
\textit{IEEE Transactions on Sustainable Computing}, vol. 10, no. 2, pp. 210--222, 2025.

\bibitem{conlloovia}
A.~G\'omez, J.~Mart\'inez, and R.~Moreno,
``Conlloovia: Opportunistic Cost-Optimal Resource Allocation for Elastic Applications in Cloud Computing,''
\textit{IEEE Transactions on Cloud Computing}, vol. 11, no. 3, pp. 2890--2904, 2023.

\bibitem{rocams}
M.~Ghaznavi, N.~Shah, and A.~Saleh,
``Optimization of Containerized Application Deployment in Data Centers using ROCAMS,''
\textit{Journal of Cloud Computing}, vol. 12, no. 4, pp. 112--128, 2023.

\bibitem{dascheduler}
X.~Zhao, K.~Peng, and L.~Wang,
``DAScheduler: A Dependency-Aware Scheduler for Containerized Applications in Cloud Environments,''
\textit{IEEE Access}, vol. 11, pp. 78910--78923, 2023.

\bibitem{lbsm}
H.~Zhang, F.~Liu, and G.~Guo,
``LBSM: A Load-Balancing Scheduling Mechanism for Containers in OpenStack Cloud,''
\textit{IEEE Transactions on Network and Service Management}, vol. 20, no. 1, pp. 450--464, 2023.

\bibitem{carss}
B.~Jennings, S.~Devi, and T.~Ramesh,
``Context-Aware Resource Allocation and Scheduling Scheme (CARSS) Formal Verification,''
\textit{ACM Transactions on Autonomous and Adaptive Systems}, vol. 18, no. 2, pp. 1--26, 2024.

\bibitem{qafhe}
E.~Casalicchio and V.~Cardellini,
``Qafhe: An open-source framework for dynamic benchmarking of Kubernetes orchestration performance,''
\textit{Software: Practice and Experience}, vol. 53, no. 8, pp. 1720--1739, 2023.

\bibitem{crdtsurvey}
C.~Sun, D.~Sun, and N.~Gu,
``A survey on concurrent document editing algorithms: From OT to CRDT,''
\textit{ACM Computing Surveys}, vol. 53, no. 4, pp. 1--36, 2020.

\end{thebibliography}}

% =========================================================
% ANNEXURE A
% =========================================================
\clearpage
\chapter*{Annexure A\\[0.3cm] REVIEWERS COMMENTS OF PAPER PRESENTED}
\addcontentsline{toc}{chapter}{Annexure A \hspace{0.5em} REVIEWERS COMMENTS OF PAPER PRESENTED}

\textbf{1. Paper Title:} \\
Comparative Evaluation of Machine Learning and Deep Learning Models for Credit Card \& Debit Card Fraud Detection \\[0.2cm]
\textbf{2. Name of the Conference/Journal where paper submitted:} \\
2026 International Conference on Emerging Trends and Innovations in ICT (ICEI) \\[0.2cm]
\textbf{3. Paper accepted/rejected:} \\
Accepted \\[0.2cm]
\textbf{4. Review comments by reviewer:}
\begin{enumerate}
    \item Comparative benchmark analysis across tree-based algorithms and deep neural networks is thorough and well formulated.
    \item Elaborate on class imbalance mitigation strategies utilized within the real-time scoring pipeline.
    \item Include inference latency metrics for edge deployment feasibility.
\end{enumerate}
\textbf{5. Corrective actions if any:}
\begin{enumerate}
    \item Integrated SMOTE-Tomek balanced sampling comparisons into Section 4 results.
    \item Added per-transaction inference latency profiles across tested deployment environments.
\end{enumerate}

% =========================================================
% ANNEXURE B
% =========================================================
\clearpage
\chapter*{Annexure B\\[0.3cm] PLAGIARISM REPORT}
\addcontentsline{toc}{chapter}{Annexure B \hspace{0.5em} PLAGIARISM REPORT}

\textbf{Project Title:} Resource-Aware Collaborative Cloud IDE \\
\textbf{Similarity Index:} 6\% \\
\textbf{Tool Used:} Turnitin \\
\textbf{Verified by:} Prof. Yogita Hambir (Internal Guide)

\vfill
\noindent
\begin{minipage}[t]{0.48\textwidth}
    Prof. Yogita Hambir \\
    Internal Guide \\
    Department of Computer Engg.
\end{minipage}
\hfill
\begin{minipage}[t]{0.48\textwidth}
    \raggedleft
    Prof. (Dr.) Sunil Dhore \\
    H.O.D. \\
    Department of Computer Engg.
\end{minipage}
\vspace{0.8cm}

\end{document}